% Modernised reading — fonds Alexandre Grothendieck, shelfmark 8, pages 1-80.
%
% This is an INTERPRETATION, not a transcription. It re-reads the pages in
% today's notation and vocabulary, states the hypotheses the manuscript leaves
% implicit, and drops the critical apparatus so that the mathematics reads
% straight through. Everything the brackets and underlines carried — a doubtful
% reading, a notation he used differently, a missing passage — moves into a
% footnote.
%
% It is held to being mathematically correct as it stands, not merely faithful
% to the page. Where the manuscript is loose or mistaken, the true statement is
% what appears here, and the footnote says what the page has. In any dispute of
% substance the transcription governs: whoever cites Grothendieck must cite the
% batch-NN.fr.tex files.
%
% Scope: the folder was read WHOLE before any of this was written — all four
% batches, pages 1-80, in one pass — and it is written as ONE file for the
% whole shelfmark. The folder is « gribouillis et calculs divers »: six
% stretches of handwritten work, separated by typescript leaves and stencilled
% versos that have nothing to do with them. The spine section names them.
%
% Notation fixed for the whole folder (see « Conventions »):
%   - phi : S_cris -> S_zar is the Frobenius morphism of topoi, epsilon the
%     canonical section, phi epsilon = F_S; phi is used for nothing else. The
%     splittings of pp. 25-26 are therefore called s, and the components of
%     Frobenius on pp. 35-41 are called F_nu (the page writes phi and phi_nu);
%   - the Galois group Gal(Kbar/K) is Gamma_K (the page writes a bold or
%     plain pi), so that pi is only ever a uniformiser;
%   - the completed algebraic closure is C_p, the Tannakian category is
%     \mathscr{C};
%   - ell^G is Illusie's co-Lie complex throughout (pp. 1-21 and 77-80 are the
%     same object, and the reading says so).
%
% Substantive departures, each footnoted where it occurs:
%   - p. 5: the sign read « != » contradicts the definition of p. 3; p. 3 kept;
%   - p. 12: « n = omega » is deduced from phi^*(n) = phi^*(omega); this needs
%     more than epsilon^* phi^* = F_S^* when S is not reduced — flagged;
%   - p. 16: H^1 of the adèles is a restricted product and includes the real
%     place; the page writes a full product over p;
%   - p. 28: the flatness condition is dA + A^A = 0, not dA + [A,A] = 0 under
%     the usual convention for [A,A];
%   - p. 30: the classifying complex's first map has the wrong sign;
%   - p. 35: the « sum of ranks = h » holds for ranks over W, not over A;
%   - p. 41: the sign hidden by an ink stain is <=, as the computation shows;
%   - p. 49: dualising t_G = pi^{-alpha} Omega^v gives omega_G = pi^alpha Omega,
%     so Omega = pi^{-alpha} omega_G, not pi^{alpha} omega_G;
%   - p. 53: the Galois action on C_sigma is semilinear;
%   - p. 62: G(Q_l) for G(Q_p), and « G = algebraic envelope » needs the
%     essential image to be closed under subobjects, not only full faithfulness;
%   - p. 80: « sans doute pas nulle si rang >= 2 » has exceptions (p = 2, d = 2).
%
% Announced and never established in the folder: the triangle (III) and its
% compatibility with F, V (p. 5, « à prouver »); the answers to the questions
% of pp. 11, 17 and 24; the « foncteur fondamental » T_p of p. 60, which the
% page uses as given.
%
% Pass: Opus 5.5 (claude-opus-5-5), 2026-10-03 — first pass, unchecked against
% the pages by a human.
\documentclass[11pt,a4paper]{article}
\input{../preamble/grothendieck}

\folder{8}
\pages{1}{80}
\dating{[vers 1967- vers 1971]}
\watermark{Édition de démonstration\\interprétation personnelle de l'œuvre}
\foldertitle{Cristaux (1970). Gribouillis et calculs divers : notes manuscrites (s.d.) — lecture modernisée du dossier entier}

\begin{document}

\section*{Résumé}

\begin{resume}
La chemise dit elle-même ce qu'elle contient : des gribouillis et des calculs
divers, autour de 1970, sous le mot « Cristaux ». Ce n'est pas un texte suivi.
Ce sont six chantiers successifs, griffonnés au dos ou entre des feuillets
d'autres textes, et qui pourtant tournent tous autour d'une même question :
comment décrire un objet arithmétique — un groupe fini en caractéristique $p$,
un groupe $p$-divisible, une variété — par de l'\emph{algèbre linéaire munie
d'un Frobenius}.

Le problème vient de loin. En caractéristique zéro, une variété complexe a sa
cohomologie, un espace vectoriel que l'on peut écrire, décomposer, comparer.
En caractéristique $p$ les outils habituels se dérobent : la cohomologie à
coefficients entiers modulo $p$ est trop pauvre, et c'est précisément le
nombre premier $p$ qui intéresse. La réponse de Grothendieck, dans ces années,
s'appelle les \emph{cristaux}. Une image aide. Un cristal est un objet qui
« pousse » de la même façon partout où on le place : connu sur l'espace
lui-même, il est connu, sans arbitraire, sur tous ses épaississements
infinitésimaux. Et sur ces cristaux agit une opération propre à la
caractéristique $p$, le Frobenius, qui élève tout à la puissance $p$. Toute
l'arithmétique se reflète alors dans la façon dont le Frobenius agit sur un
module : le même groupe fini, qui en lui-même est un objet très rigide et
difficile à saisir, devient un module libre muni de deux flèches $F$ et $V$,
dont le composé est la multiplication par $p$.

Le dossier s'ouvre sur la question la plus fine : un groupe fini tué par $p$,
dont le « cristal de Dieudonné » vit sur un site compliqué (le site
cristallin), peut-il être décrit par des données vivant sur l'espace ordinaire,
le site de Zariski ? L'outil est une remarque simple — en caractéristique $p$,
le Frobenius d'un épaississement à puissances divisées se factorise par
l'espace lui-même — qui fournit un pont dans un seul sens, et la question est
de savoir ce que l'on perd en le traversant. Viennent ensuite des calculs
explicites, en coordonnées, de modules de Dieudonné : pour les groupes
« ordinaires », les plus simples qui soient, puis pour les groupes sur
lesquels opère l'anneau des entiers d'un corps $p$-adique. Ces derniers
mènent aux groupes de Lubin et Tate, qui fabriquent la théorie du corps de
classes local, et à un théorème de Tate qui décompose leur module de Tate
comme on décompose la cohomologie d'une variété complexe en types de Hodge.

À partir de là le regard change d'échelle. Les dernières pages regardent
toutes ces structures d'un coup — le Frobenius cristallin, la filtration de
Hodge, l'action du groupe de Galois, la comparaison avec la cohomologie
complexe — comme les \emph{réalisations} d'un même objet, et cherchent le
groupe dont elles sont les représentations : le langage est celui des
catégories tannakiennes, et l'objet est le motif. Une page dresse la liste des
conjectures en jeu, Hodge et Tate en tête. Le dossier se ferme sur un calcul
modeste, qui rejoint la première page : le complexe qui mesure la géométrie
infinitésimale d'un groupe de type multiplicatif, et l'obstruction, cachée
dedans, à relever une représentation modulo $p$ en une représentation modulo
$p^{2}$.

On y voit une manière de travailler. Grothendieck pose ses données par
lettres, a), b), c), en se demandant chaque fois ce qu'elles déterminent et ce
qu'elles laissent libre ; il écrit « à prouver ! » à côté d'un énoncé dont il a
besoin, « Question » devant une équivalence qu'il espère, et un point
d'interrogation devant une conjecture qu'il se propose de « revoir ».
Plusieurs des questions posées ici ont reçu depuis des réponses, sous d'autres
noms : la théorie de Dieudonné cristalline, les représentations cristallines,
les périodes $p$-adiques. Les pages ne le savent pas, et la lecture qui suit
s'efforce de ne pas le leur faire dire.

\keywords{Dieudonné crystal, crystalline site, Frobenius descent, co-Lie complex, finite flat group scheme killed by p, truncated Barsotti-Tate group, ordinary Barsotti-Tate group, F-crystal with connection, Frobenius lift, formal A-module, Lubin-Tate group, Grothendieck-Messing deformation theory, Hodge-Tate decomposition, Lubin-Tate character, Tannakian category, filtered F-isocrystal, mysterious functor, motivic Galois group, period torsor, Hodge conjecture, Tate conjecture, group of multiplicative type, lifting obstruction}
\end{resume}

\pagerange{1}{1}
\section*{Le fil du dossier, et les conventions}

\subsection*{Les stations}

Le dossier compte quatre-vingts pages, dont un peu plus de la moitié sont de
la main de Grothendieck. Le reste se partage entre des feuillets de
tapuscrits portant des corrections à la main, qui n'ont pas de rapport avec
les notes voisines et sont décrits là où ils tombent, et des versos
dactylographiés ou ronéotypés sans aucune intervention (un texte anglais sur
les points doubles pseudo-rationnels, une « Tribu » de Bourbaki, deux cours
polycopiés sur les variétés $C^{r}$ et sur l'intégration), que la lecture
ignore. Les pages de sa main forment six chantiers :

\begin{enumerate}
\item \emph{Pages 1 à 21.} Descente du cristal de Dieudonné d'un groupe fini
tué par $p$ au site de Zariski, au moyen des complexes co-Lie de $G$ et de son
dual. Le programme est posé trois fois (pages 3--4, 9, 15--17), calculé dans
le cas des groupes de Barsotti--Tate tronqués d'échelon 1 (pages 12--13), et
ramené pour finir à un problème de descente dont l'indétermination est
mesurée par des $\operatorname{Ext}$ relatifs (page 19). Une page barrée, la
page 16, traite d'autre chose.
\item \emph{Pages 24 à 30.} Le module de Dieudonné d'un groupe de
Barsotti--Tate ordinaire sur une base lisse, et plus généralement d'un
$F$-cristal à connexion, calculé en coordonnées après choix d'un relèvement
du Frobenius.
\item \emph{Pages 33 à 55.} Les groupes de Barsotti--Tate sur
$\mathbb{F}_{q}$ sur lesquels opère l'anneau des entiers $A$ d'un corps
$p$-adique : classification par le module de Dieudonné, caractérisation de
la dimension 1, relèvement en groupe de Lubin--Tate, et lecture du théorème de
Tate sur ce relèvement.
\item \emph{Pages 57 à 59.} Trois feuilles de papier quadrillé : des lignes
polygonales et des comparaisons de fractions.
\item \emph{Pages 60 à 74.} Le groupe de Galois tannakien des $F$-cristaux
filtrés à isogénie près, puis la même liste de structures pour les motifs sur
un ouvert de $\operatorname{Spec}\mathbb{Z}$, et une liste de conjectures.
\item \emph{Pages 77 à 80.} Le complexe co-Lie d'un groupe de type
multiplicatif, et la classe d'obstruction qu'il contient.
\end{enumerate}

Le premier et le dernier chantier parlent du même objet, le complexe co-Lie
$\ell^{G}$ ; le deuxième et le troisième calculent des modules de Dieudonné ;
le troisième et le cinquième parlent tous deux de comparaison entre le module
de Tate et la filtration de Hodge. Rien dans le dossier ne dit dans quel
ordre ces pages ont été écrites, et l'on ne prétend pas qu'elles forment un
argument : ce sont des stations, et la lecture les prend dans l'ordre des
pages.

\subsection*{Conventions, valables pour tout le dossier}

$p$ est un nombre premier. Pour un schéma $S$ de caractéristique $p$, $F_{S}$
est son Frobenius absolu et, pour un $\mathcal{O}_{S}$-module $E$,
$E^{(p)} = F_{S}^{*}E$. On écrit $S_{\mathrm{zar}}$ pour le topos de Zariski de
$S$ et $S_{\mathrm{cris}}$ pour le topos cristallin de $S$ sur
$\mathbb{F}_{p}$ (sur $\mathbb{Z}_{p}$ quand on le précise).

\emph{Le morphisme $\varphi$.} La lettre $\varphi$ désigne dans tout le
document le morphisme de topos
\[
\varphi\colon S_{\mathrm{cris}} \longrightarrow S_{\mathrm{zar}},
\qquad \varphi\,\varepsilon = F_{S},
\]
où $\varepsilon\colon S_{\mathrm{zar}} \to S_{\mathrm{cris}}$ est le
morphisme canonique. Il existe pour la raison que la page 19 dessine : si
$(U, T, \delta)$ est un épaississement à puissances divisées en
caractéristique $p$, d'idéal $J$, alors $x^{p} = p!\,\gamma_{p}(x) = 0$ pour
$x \in J$, donc le Frobenius de $T$ se factorise par $U$, et l'on pose
$\varphi^{*}(E)_{T} = (T \to U \to S)^{*}E$. Pour un cristal $\mathcal{M}$
sur $S_{\mathrm{cris}}$, $\varphi^{*}\varepsilon^{*}\mathcal{M} =
\mathcal{M}^{(p)}$ est sa tordue par Frobenius, et pour un module $E$ sur
$S_{\mathrm{zar}}$, $\varepsilon^{*}\varphi^{*}E = E^{(p)}$.\footnote{La page
1 pose seulement $\varphi\varepsilon = F_{S}$ (« frobenius absolu ») et la
page 12 utilise $\varphi^{*}(M) = \mathcal{M}^{(p)}$ ; la justification par
la factorisation du Frobenius est ce que le diagramme $J \hookrightarrow B
\to A$, $\lambda \mapsto \lambda^{p}$, de la page 19 esquisse. Le diagramme
est très retouché, et la lecture qu'on en donne ici est la nôtre. La lettre
$\varphi$ sert aussi, sur la page, aux scindages des pages 25--26 et aux
composantes du Frobenius des pages 35--41 : on les renomme $s$ et $F_{\nu}$
pour lui garder un seul sens.}

\emph{Le cristal et sa valeur.} $\mathcal{M}$ (lettre ronde) est un cristal,
$M = \varepsilon^{*}\mathcal{M}$ (lettre droite) sa valeur sur
$S_{\mathrm{zar}}$. Pour un groupe fini localement libre $G$,
$\mathbb{D}^{*}(G)$ est son cristal de Dieudonné (contravariant), muni de
$F\colon \mathbb{D}^{*}(G)^{(p)} \to \mathbb{D}^{*}(G)$ et
$V\colon \mathbb{D}^{*}(G) \to \mathbb{D}^{*}(G)^{(p)}$, et $G^{*}$ est le
dual de Cartier de $G$.\footnote{Le cristal de Dieudonné d'un groupe fini
localement libre arbitraire n'a été construit en toute généralité que par
Berthelot, Breen et Messing (\emph{Théorie de Dieudonné cristalline II},
1982) ; pour les groupes de Barsotti--Tate, Grothendieck l'avait annoncé au
congrès de Nice (1970), et Mazur et Messing l'ont rédigé (1974). Les pages 3
et 4 le posent comme donné « conjecturalement ».}

\emph{Le complexe co-Lie.} $\ell^{G}$ est le complexe co-Lie de $G$ au sens
d'Illusie, de tor-amplitude $[-1, 0]$ pour $G$ fini localement libre, avec
\[
\omega_{G} = H^{0}(\ell^{G}), \qquad n_{G} = H^{-1}(\ell^{G}), \qquad
t_{G} = \omega_{G}^{\vee} = \operatorname{Lie}(G).
\]
La page écrit tantôt $\ell^{G}_{\bullet}$, tantôt $\ell_{\bullet}$ pour
$\ell^{G}$ et $\ell'_{\bullet}$ pour $\ell^{G^{*}}$ ; le « $\check{\ ~}$ »
désigne le dual $R\mathcal{H}om(-, \mathcal{O})$.\footnote{L'identification
des « complexes parfaits $\ell_{\bullet}$, $\ell'_{\bullet}$ d'amplitude
$[-1,0]$ » des pages 3--4 avec les complexes co-Lie de $G$ et $G^{*}$ est
appuyée par trois choses : la notation $\ell^{G}_{\bullet}$, $\ell^{G^{*}}_{\bullet}$
de la page 4 ; le calcul de la page 12, où $H^{0}(\ell) = \omega$ et
$H^{-1}(\ell) = n$ ; et les pages 77--80, où $\ell^{G}_{\bullet}$ est
explicitement le complexe co-Lie, calculé par le triangle de transitivité. Le
nom de « complexe co-Lie » est celui d'Illusie (\emph{Complexe cotangent et
déformations II}, 1972) ; la page 77 attribue le triangle à Mazur et Roberts.}

\emph{Les corps locaux} (pages 33 à 74). $K$ est une extension finie de
$\mathbb{Q}_{p}$, d'anneau d'entiers $A$, d'uniformisante $\pi$, d'indice de
ramification $e$, de corps résiduel $k = \mathbb{F}_{q}$, $q = p^{f}$, et
$h = ef = [K : \mathbb{Q}_{p}]$ ; $W = W(k)$, $K_{0} = W[1/p]$, $\sigma$ est le
Frobenius de $W$. $\Gamma_{K} = \operatorname{Gal}(\bar K/K)$, et
$\mathbb{C}_{p}$ est le complété de $\bar K$.\footnote{La page écrit $C$ pour
$\mathbb{C}_{p}$ et $\pi$ (en gras pages 51--55, maigre pages 60--70) pour le
groupe de Galois. On réserve $\pi$ à l'uniformisante.}

\pagerange{1}{5}
\section*{Le cristal d'un groupe tué par $p$ et sa descente (pages 1 à 5)}

\subsection*{Les données (pages 3 et 4)}

Soit $G$ un schéma en groupes commutatif fini localement libre, tué par $p$,
sur une base $S$ de caractéristique $p$. Le programme, posé page 4 comme une
liste de structures qu'on « associe (conjecturalement) » à $G$ et repris page 3
pour un cristal abstrait, est le suivant.

a) Un cristal $\mathcal{M} = \mathbb{D}^{*}(G)$ sur $S_{\mathrm{cris}}$,
localement libre de type fini, fonctoriel en $G$, avec
\[
\mathcal{M}^{(p)} \xrightarrow{\ F\ } \mathcal{M} \xrightarrow{\ V\ }
\mathcal{M}^{(p)}, \qquad FV = 0, \quad VF = 0.
\]

b) Un complexe parfait $\Delta = \Delta(\mathcal{M})$ sur $S_{\mathrm{cris}}$,
de tor-amplitude dans un intervalle de longueur 1, dont les deux faisceaux
de cohomologie sont isomorphes à $\mathcal{M}$, muni de $F_{\Delta}$,
$V_{\Delta}$ avec $F_{\Delta}V_{\Delta} = V_{\Delta}F_{\Delta} = 0$, et deux
triangles distingués
\[
(\mathrm{I}) \quad [\mathcal{M}^{(p)} \xrightarrow{F} \mathcal{M}]
\xrightarrow{\ i_{\mathrm{I}}\ } \Delta^{(p)}
\xrightarrow{\ \pi_{\mathrm{I}}\ } [\mathcal{M} \xrightarrow{V} \mathcal{M}^{(p)}]
\xrightarrow{\ \partial_{\mathrm{I}}\ } ,
\]
\[
(\mathrm{II}) \quad [\mathcal{M} \xrightarrow{V} \mathcal{M}^{(p)}]
\xrightarrow{\ i_{\mathrm{II}}\ } \Delta
\xrightarrow{\ \pi_{\mathrm{II}}\ } [\mathcal{M}^{(p)} \xrightarrow{F} \mathcal{M}]
\xrightarrow{\ \partial_{\mathrm{II}}\ } ,
\]
avec $\Delta^{(p)} \simeq \varphi^{*}(\varepsilon^{*}\Delta)$.\footnote{La
page 3 écrit l'amplitude « $\subset [0,1]$ », la page 9 « $\subset [-1,0]$ » ;
c'est une affaire de décalage, et la page 9 est la bonne pour la construction
qu'on donne plus bas, où $\Delta = \mathcal{D} \otimes^{L} \mathbb{Z}/p$. La
flèche horizontale de (I) est notée sur la page tantôt $\partial_{\mathrm{I}}$,
tantôt $i_{\mathrm{I}}$ ; on garde $i_{\mathrm{I}}$.}

c) Sur $S_{\mathrm{zar}}$, les complexes co-Lie $\ell = \ell^{G}$ et
$\ell' = \ell^{G^{*}}$, de tor-amplitude $[-1,0]$, et un accouplement
$(\ell^{G^{*}})^{\vee} \to \ell^{G}$, c'est-à-dire une flèche
$\mathcal{O}_{S} \to \ell^{G^{*}} \otimes^{L} \ell^{G}$.\footnote{Page 4, la
rubrique « b) homomorphismes » reste sans suite, et la rubrique « e) Un
triangle exact » s'arrête au bas de la page. La flèche $\pi$ de la page porte
le nom de la projection d'un triangle ; on l'appelle ici accouplement.}

d) Un triangle distingué sur $S_{\mathrm{zar}}$, posant
$\Delta_{\mathrm{zar}} = \varepsilon^{*}\Delta$,
\[
(\mathrm{III}) \quad \ell[-1] \xrightarrow{\ i_{0}\ } \Delta_{\mathrm{zar}}
\xrightarrow{\ \pi_{0}\ } \check\ell' \xrightarrow{\ \partial_{0}\ } ,
\]
et des isomorphismes $\varphi^{*}(\ell[-1]) \simeq [\mathcal{M}^{(p)}
\xrightarrow{F} \mathcal{M}]$, $\varphi^{*}(\check\ell') \simeq
[\mathcal{M} \xrightarrow{V} \mathcal{M}^{(p)}]$ qui identifient
$\varphi^{*}(\mathrm{III})$ au triangle (I).

C'est là tout le sens du programme. Le triangle (I) est une « filtration »
du complexe $\varphi^{*}\Delta_{\mathrm{zar}} = \Delta^{(p)}$ sur le site
cristallin ; (III) en est une \emph{descente} au site de Zariski, et ce qui
descend, ce sont les complexes co-Lie de $G$ et de $G^{*}$. Si $M =
\varepsilon^{*}\mathcal{M}$ (un module localement libre muni d'une connexion
intégrable), on en tire en appliquant $\varepsilon^{*}$
\[
\ell[-1]^{(p)} \simeq [M^{(p)} \xrightarrow{F_{M}} M], \qquad
\check\ell'^{(p)} \simeq [M \xrightarrow{V_{M}} M^{(p)}],
\]
ce que la page encadre d'un point d'exclamation, avec en marge « comp. avec
connexion ».

\subsection*{Le cas où $F$ est nul (page 1)}

La page 1, écrite avant ou à côté de ce programme, en calcule un cas. Si $G$
est de hauteur $\le 1$ — le noyau du Frobenius d'un groupe lisse $H$ —, son
cristal est $\varphi^{*}(\omega)$ avec $\omega = \omega_{G}$, le Frobenius
du cristal est nul, et $V$ est $\varphi^{*}$ de la flèche
\[
C\colon \omega \longrightarrow \omega^{(p)},
\]
transposée de l'application puissance $p$-ième de l'algèbre de Lie.\footnote{L'encadré
de la page porte aussi « $\Pi\colon \nu^{(p)} \to \nu$ », de lecture
douteuse, sans emploi dans la suite.} Le complexe co-Lie se calcule alors à
la main : $G = \operatorname{Ker}(F_{H})$ donne
$\ell^{G} \simeq [\omega^{(p)} \xrightarrow{0} \omega]$, la différentielle du
Frobenius étant nulle, et le triangle (III) devient
\[
[\omega^{(p)} \xrightarrow{0} \omega]
\xrightarrow{(\mathrm{id},\,C)}
[\omega^{(p)} \xrightarrow{0} \omega^{(p)}]
\xrightarrow{(0,\,\mathrm{id})}
[\omega \xrightarrow{C} \omega^{(p)}]
\xrightarrow{(\mathrm{id},\,0)} ,
\]
où le terme du milieu est $\Delta_{\mathrm{zar}}$, à cohomologie $\omega^{(p)}
= M$ en chaque degré.\footnote{On a vérifié que ce triangle est distingué :
le cône de $(\mathrm{id}, C)$ est quasi-isomorphe à
$[\omega \xrightarrow{C} \omega^{(p)}]$, le noyau de la comparaison étant le
complexe acyclique $[\omega^{(p)} \xrightarrow{\mathrm{id}} \omega^{(p)}]$.
La page dessine un second triangle, tordu par $p$, où les décalages des
sommets ne s'accordent pas entre eux (un $[1]$ là où l'amplitude demande un
$[-1]$) ; on ne le reproduit pas. Le sommet du second triangle de gauche est
surchargé.} Le haut de la page note aussi la suite exacte
$0 \to {}_{F}\mathcal{M}^{(p)} \to \mathcal{M} \to {}_{V}\mathcal{M}^{(p)}
\to$, inachevée.

\subsection*{Compatibilités avec $F$ et $V$ (pages 3 et 5)}

Sur $\ell$ et $\ell'$, les structures sont $V_{\ell}\colon \ell \to
\ell^{(p)}$, $F_{\ell} = 0$, et de même pour $\ell'$. La page 3 les
\emph{définit} à partir des triangles :
\[
F_{\check\ell'} := \pi_{0}\,\varepsilon^{*}(i_{\mathrm{II}}) = {}^{t}V_{\ell'},
\qquad
V_{\ell}[-1] := \varepsilon^{*}(\pi_{\mathrm{II}})\, i_{0}.
\]
La page 5 annonce deux énoncés : 1) le triangle (III) est compatible aux
structures $(F, V)$ — « \emph{à prouver} ! » ; 2) l'isomorphisme
$\varphi^{*}(\mathrm{III}) \simeq (\mathrm{I})$ l'est aussi. Elle vérifie
que $i_{0}$ est compatible à $V$ et $\pi_{0}$ à $F$ (le reste « de même »),
que $\partial_{0}$ est compatible à $F$ parce que
$\partial_{0}\,\pi_{0}\,\varepsilon^{*}(i_{\mathrm{II}}) = 0$ (deux flèches
consécutives d'un triangle), et pour 2) le calcul
\[
\varphi^{*}(V_{\ell}[-1]) = \varphi^{*}\varepsilon^{*}(\pi_{\mathrm{II}})\,
\varphi^{*}(i_{0}) = \pi_{\mathrm{II}}^{(p)}\, i_{\mathrm{I}}
= V_{[\mathcal{M}^{(p)} \to \mathcal{M}]},
\]
« et l'autre relation se prouve de même !! ».\footnote{Au milieu de la page,
une ligne se lit « $\varepsilon^{*}(\pi_{\mathrm{II}})\,i_{0} \ne
V_{\ell}[-1]$ » — le signe est une égalité barrée, de lecture douteuse. Elle
contredit la définition de la page 3, qu'on suit ; il s'agit vraisemblablement
d'un signe de définition et non d'une négation.} Le point 1), en tant
qu'existence de (III), n'est établi nulle part dans le dossier.

\pagerange{6}{14}
\subsection*{Un tapuscrit intercalé (pages 6, 8, 10 et 14)}

Quatre feuillets d'un tapuscrit français, numérotés 32, 33, 35 et 37.1,
s'intercalent entre ces notes. Ils portent la fin de la démonstration d'un
théorème de changement de base propre pour les faisceaux de torsion (un
théorème (5.1) à trois volets (i), (ii), (iii) pour des triples
$(g, f, F)$), le Lemme (8.3) — si $\pi\colon \bar X \to X$ est propre et
surjectif au-dessus de $S$, les énoncés (5.1) pour $\bar f$ et pour $\pi$
entraînent ceux pour $f$ — et le Lemme (8.4), qui ramène (5.1) au cas $f$
projectif. Les corrections à la main sont de rédaction : « s'ensuit »
remplacé par « résulte », « applicable » par « pertinent », l'énoncé (i) du
Lemme (8.3) biffé et remplacé par « Assertion vide », un « un peu bref, faire
l'argument » en marge, des renvois complétés.\footnote{La numérotation
(5.1), (8.3), (8.4) et les renvois « IX 1.2 » évoquent la rédaction des
exposés de SGA 4 sur le changement de base propre ; on ne l'a pas vérifié, et
la main des corrections n'est pas identifiée.} Ils n'ont pas de rapport avec
les cristaux.

\pagerange{7}{9}
\section*{D'où viennent les deux triangles (pages 7 et 9)}

\subsection*{$F$ et $V$ se lisent sur les triangles (page 7)}

La page 7 récrit les deux triangles avec des lettres abstraites :
$A \to \mathfrak{X}^{(p)} \to B \to$ pour (I) et $B \to \mathfrak{X} \to A
\to$ pour (II). Elle observe que le Verschiebung de $\mathfrak{X}$ est le
composé
\[
V_{\mathfrak{X}} = i_{\mathrm{I}}\,\pi_{\mathrm{II}}\colon
\mathfrak{X} \to A \to \mathfrak{X}^{(p)},
\]
et, de même, $F_{\mathfrak{X}} = i_{\mathrm{II}}\,\pi_{\mathrm{I}}\colon
\mathfrak{X}^{(p)} \to B \to \mathfrak{X}$.\footnote{La page n'écrit
explicitement que $V_{\mathfrak{X}}$ (le premier membre de l'encadré est
surchargé, lu $\pi_{\mathrm{I}}$ ou $\pi^{(p)}_{\mathrm{II}}\,i_{\mathrm{I}}$) ;
l'expression de $F_{\mathfrak{X}}$ est notre complément, par symétrie.} Les
relations $F_{\mathfrak{X}}V_{\mathfrak{X}} = 0$ et
$V_{\mathfrak{X}}F_{\mathfrak{X}} = 0$ sont alors automatiques, puisque
$\pi_{\mathrm{I}}\,i_{\mathrm{I}} = 0$ et $\pi_{\mathrm{II}}\,i_{\mathrm{II}} = 0$.
Il se demande encore si les composés des cobords
$\partial_{\mathrm{I}}\partial_{\mathrm{II}}$ et
$\partial_{\mathrm{II}}\partial_{\mathrm{I}}$ sont nuls (« nul ? », entouré
deux fois), et note que $\varepsilon^{*}$ envoie (I) sur la tordue
$(\mathrm{II})^{(p)}$.

\subsection*{Le cristal sur $\mathbb{Z}_{p}$ (page 9)}

La page 9 reprend tout avec une donnée plus naturelle, et c'est elle qui
explique les triangles. On part d'un cristal $\mathcal{D}$ sur
$S_{\mathrm{cris}}/\mathbb{Z}_{p}$, de présentation finie et de tor-amplitude
$[-1, 0]$, avec $F$ et $V$ — par exemple le cristal de Dieudonné de $G$, qui
est tué par $p$ si $G$ l'est. On pose
\[
\Delta = \mathcal{D} \otimes^{L} \mathbb{Z}/p
\]
sur $S_{\mathrm{cris}}/\mathbb{F}_{p}$. Si $\mathcal{D}$ est tué par $p$,
les deux faisceaux de cohomologie de $\Delta$ sont $\mathcal{D}[p] = \mathcal{D}$
et $\mathcal{D}/p = \mathcal{D}$ : c'est exactement la propriété demandée
page 3. Et les triangles sont ceux de l'axiome de l'octaèdre appliqué aux
deux factorisations de $p$ :
\[
p = V F \ \text{ sur } \mathcal{D}^{(p)} \ \Longrightarrow\
\mathcal{D}\overset{L}{\otimes}/F \to \mathcal{D}^{(p)}\overset{L}{\otimes}/p
\to \mathcal{D}\overset{L}{\otimes}/V \to ,
\]
\[
p = F V \ \text{ sur } \mathcal{D} \ \Longrightarrow\
\mathcal{D}\overset{L}{\otimes}/V \to \mathcal{D}\overset{L}{\otimes}/p
\to \mathcal{D}\overset{L}{\otimes}/F \to ,
\]
où « $\otimes^{L}/u$ » désigne le cône de $u$. Ce sont (I) et (II), avec
$\mathcal{D}^{(p)}\otimes^{L}/p \simeq \varphi^{*}(\Delta_{\mathrm{zar}})$.\footnote{La
page écrit les triangles et l'identification
$\mathcal{M}^{(p)}\otimes^{L}/p \simeq \varphi^{*}(\Delta^{*})$, sans nommer
l'octaèdre ; l'explication est la nôtre. Sous le second triangle, une
inscription se lit « $\simeq^{\mathrm{dbs}}$ $\otimes^{L}\Delta^{*}(\mathcal{M})$ »,
de lecture douteuse.}

Le point b) devient alors, sans conjecture : \emph{descendre} à
$S_{\mathrm{zar}}$, via $\varphi$, la filtration (I) de
$\varphi^{*}(\Delta_{\mathrm{zar}})$, c'est-à-dire trouver le triangle (III).
Et la page va un cran plus loin : sur un épaississement à puissances divisées
$S'$ de $S$, on dispose de $\Delta(S') = \mathcal{D} \otimes^{L}
\mathcal{O}_{S'}$, qui prolonge $\Delta$, et il faut se donner un prolongement
$\bar\ell[-1] \to \Delta(S') \to \check{\bar\ell}'$ du triangle
(III).\footnote{« Épaississement à puissances divisées » et la parenthèse qui
suit sont des lectures incertaines ; la parenthèse paraît fixer l'idéal de
l'épaississement en fonction de $p$ et de l'idéal de $S$.} C'est l'esquisse
d'une théorie des déformations par les cristaux : déformer $G$, c'est
prolonger la filtration.

\pagerange{11}{13}
\section*{Questions, et le cas des groupes tronqués d'échelon 1 (pages 11 à 13)}

\subsection*{Quatre questions (page 11)}

a) Les catégories de structures a), b) définissent-elles des \emph{champs}
sur $S_{\mathrm{zar}}$ ?

b) Le foncteur qui oublie la descente et ne garde que le cristal
$\mathcal{M}$ sur $S_{\mathrm{cris}}/\mathbb{F}_{p}$ est-il \emph{fidèle} ?\footnote{« Restriction
partiel » est une lecture incertaine. Une note de marge, en partie biffée et
encadrée, lit par fragments une suite exacte d'homologie se terminant par
« $0 \to t^{*} \to \omega$ nul, et $\omega^{*}$ et $t^{*}$ loc. libres ».}

c) Le foncteur naturel qui va des modules localement libres $\omega$ sur
$S_{\mathrm{zar}}$, munis de $F_{\omega}$, $V_{\omega}$, vers les structures
a), b) est fidèle ; est-il pleinement fidèle ? L'est-il du moins sur la
sous-catégorie $F_{\omega} = 0$, ou sur celle où
$\operatorname{Ker} F = \operatorname{Im} V$ ?

d) Peut-on à $(\omega, F_{\omega}, V_{\omega})$ associer un schéma en groupes
fini localement libre sur $S$ ?\footnote{Pour $F_{\omega} = 0$ et $G$
commutatif, la question d) a une réponse classique : les groupes de hauteur
$\le 1$ correspondent aux $p$-algèbres de Lie localement libres (SGA 3,
exposé VII$_{\mathrm{A}}$), et pour un groupe commutatif la donnée est celle
de $\omega$ et de la transposée $V_{\omega}$ de l'application puissance
$p$-ième. La page ne le dit pas.}

\subsection*{Les groupes de Barsotti--Tate tronqués d'échelon 1 (pages 12 et 13)}

Supposons $G$ de Barsotti--Tate tronqué d'échelon 1, c'est-à-dire que
\[
\mathcal{M}^{(p)} \xrightarrow{F} \mathcal{M} \xrightarrow{V}
\mathcal{M}^{(p)} \xrightarrow{F} \mathcal{M}
\]
est exacte, ce qui équivaut à l'exactitude de la suite correspondante pour
$M$. Prenons la cohomologie du triangle (III). En degré 1 on obtient, avec
$\omega = H^{1}(\ell[-1]) = H^{0}(\ell)$ et $\nu' = H^{1}(\check\ell')$,
tous deux localement libres, une filtration
\[
0 \to \omega \to M \to \nu' \to 0
\]
telle que, dans $\varphi^{*}(M) = \mathcal{M}^{(p)}$, le sous-cristal
$\varphi^{*}(\omega)$ soit $\operatorname{Im}(V\colon \mathcal{M} \to
\mathcal{M}^{(p)})$. En degré 0, avec $n = H^{-1}(\ell)$ et $t' =
H^{0}(\check\ell')$,
\[
0 \to n \to M \to t' \to 0, \qquad
\varphi^{*}(n) = \operatorname{Ker}(F\colon \mathcal{M}^{(p)} \to \mathcal{M}).
\]
Comme $\operatorname{Ker} F = \operatorname{Im} V$, on a
$\varphi^{*}(n) = \varphi^{*}(\omega)$, et la page conclut
\[
n = \omega, \qquad t' = \nu' .
\]
On retrouve deux faits connus : la tordue par Frobenius de la filtration de
Hodge $0 \to \omega_{G} \to M \to t_{G^{*}} \to 0$ est
$\operatorname{Ker} F = \operatorname{Im} V$, et $n_{G} \simeq \omega_{G}$ —
ce qu'on voit aussi en écrivant localement $G = X[p]$ pour un groupe de
Barsotti--Tate $X$, d'où $\ell^{G} \simeq [\omega_{X} \xrightarrow{p}
\omega_{X}]$, de différentielle nulle en caractéristique
$p$.\footnote{La page écrit
en marge « Cela détermine $\omega$ de façon \emph{unique} ! ». La déduction
de $n = \omega$ à partir de $\varphi^{*}(n) = \varphi^{*}(\omega)$ demande que
$\varphi^{*}$ distingue les sous-modules localement facteurs directs de $M$.
Ce n'est pas une conséquence de $\varepsilon^{*}\varphi^{*} = F_{S}^{*}$
seule : sur une base non réduite, $F_{S}^{*}$ ne distingue pas deux tels
sous-modules qui coïncident modulo le nilradical. Il faut utiliser
$\varphi^{*}$ sur tous les épaississements, et la page ne le justifie pas.
Sur $S$ réduit, l'argument est complet tel quel.}

Reste à savoir si $\mathcal{M}$, $F$, $V$ suffisent à reconstituer toute la
donnée b). On peut déjà construire
\[
V_{\omega}\colon \omega \to \omega^{(p)}, \qquad
F_{t'}\colon t'^{(p)} \to t',
\]
le premier parce que $\omega^{(p)} = \operatorname{Im}(V_{M})$, en restreignant
$V_{M}$ à $\omega$ ; le second parce que $t'^{(p)} =
\operatorname{Coker}(V_{M})$ et que $F_{M}$, nul sur l'image de $V_{M}$, se
factorise en une injection $t'^{(p)} \hookrightarrow M$, qu'on compose avec
$M \to t'$. Le dual $\omega' = \check t'$ reçoit ainsi $V_{\omega'} =
{}^{t}F_{t'}$. La page 13 dessine les carrés correspondants, et une esquisse
inachevée où les deux constructions se rejoignent sur la suite
$0 \to \omega \to M \to t' \to 0$.\footnote{L'esquisse de la page 13 n'a pas
de pointes de flèches sur sa ligne horizontale, et la place de $\omega$ n'y
est qu'approchée.}

\pagerange{15}{17}
\section*{Le programme en quatre points (pages 15 et 17)}

La page 15 reformule la donnée pour un cristal $\mathcal{M}$ de présentation
finie sur $S_{\mathrm{cris}}/\mathbb{F}_{p}$ avec $F$, $V$ :

1° les triangles $\mathrm{I}(\mathcal{M})\colon L(\mathcal{M}) \to
\Delta(\mathcal{M})^{(p)} \to L'(\mathcal{M}) \to$ et
$\mathrm{II}(\mathcal{M})$, avec $\varepsilon^{*}\mathcal{M} = M$ ;

2° un triangle $\mathcal{T}\colon \mathcal{T}_{0} \to \mathcal{T}_{1} \to
\mathcal{T}_{2} \to$ de complexes parfaits de tor-amplitude dans $[-1,0]$ sur
$S_{\mathrm{zar}}$, peut-être muni de $F$ et $V$ (« F-V-triangles ? ») ;

3° un isomorphisme de triangles $\alpha\colon \varphi^{*}(\mathcal{T}) \simeq
\mathrm{I}(\mathcal{M})$ ;

4° un isomorphisme $\beta\colon \mathcal{T}_{1} \simeq
\varepsilon^{*}\Delta(\mathcal{M})$ tel que $\varphi^{*}(\beta)$ et
$\alpha_{1}$ coïncident, modulo l'identification
$\varphi^{*}\varepsilon^{*}\Delta(\mathcal{M}) = \Delta(\mathcal{M})^{(p)}$.\footnote{La
page note le sommet du triangle $\mathcal{T}_{1}$, comme le coin de droite ;
l'un des deux porte vraisemblablement un autre indice. On les nomme
$\mathcal{T}_{0}$, $\mathcal{T}_{1}$, $\mathcal{T}_{2}$ en suivant le coin
qui doit être $\Delta$. En tête de page, isolés, une suite
$0 \to G \to G' \to G'' \to 0$ et une inclusion de sites cristallins
$X_{\mathrm{cris}/0} \hookrightarrow X_{\mathrm{cris}/n}$.}

La condition 4° est le point délicat : $\mathcal{T}_{1}$, sur le site de
Zariski, doit \emph{être} la valeur de $\Delta$, et pas seulement en devenir
une après $\varphi^{*}$.

La page 17 en donne l'exemple qui guide tout. À $\omega$ sur
$S_{\mathrm{zar}}$ muni de $F_{\omega}$, $V_{\omega}$, on associe :
$\mathcal{M} = \varphi^{*}(\omega)$, $F_{\mathcal{M}} = \varphi^{*}(F_{\omega})$,
$V_{\mathcal{M}} = \varphi^{*}(V_{\omega})$, donc $M = \omega^{(p)}$ ; le
triangle
\[
\mathcal{T}(\omega)\colon\quad
[\omega^{(p)} \xrightarrow{F_{\omega}} \omega] \longrightarrow
[\omega^{(p)} \xrightarrow{0} \omega^{(p)}] \longrightarrow
[\omega \xrightarrow{V_{\omega}} \omega^{(p)}] \longrightarrow ,
\]
dont le terme du milieu est $\Delta(\omega) = \varepsilon^{*}\Delta(\mathcal{M})$ ;
des isomorphismes canoniques $\alpha_{\omega}\colon \varphi^{*}\mathcal{T}(\omega)
\simeq \mathcal{T}(\varphi^{*}\omega)$ et $\beta_{\omega}$ ; et pour 4°,
l'identité.\footnote{Le coin gauche du triangle est lu « $\omega^{V}$ »,
peut-être $\omega^{(p)}$ ; c'est cette dernière lecture que demande la
flèche $F_{\omega}$, et qu'on suit.}

\emph{Question.} Le foncteur
\[
(\omega, F_{\omega}, V_{\omega}) \longmapsto
\bigl(\varphi^{*}\omega,\ \varphi^{*}F_{\omega},\ \varphi^{*}V_{\omega},\
\mathcal{T}(\omega),\ \alpha_{\omega},\ \beta_{\omega}\bigr)
\]
est-il pleinement fidèle ? L'est-il du moins sur la sous-catégorie
$F_{\omega} = 0$ ? Un morphisme de but est une paire $(u, v)$, $u\colon
\varphi^{*}\omega \to \varphi^{*}\omega'$ compatible à $F$ et $V$,
$v\colon \mathcal{T}(\omega) \to \mathcal{T}(\omega')$ un morphisme de
triangles, rendant commutatifs les carrés formés avec $\alpha$ et $\beta$ ;
la page note que $v_{1}$ est alors déterminé comme $\varepsilon^{*}(u)$, et
$\varphi^{*}(v)$ comme $\mathcal{T}(u)$. La question reste ouverte dans le
dossier.

\pagerange{16}{16}
\subsection*{Une page barrée : rigidité d'une donnée adélique (page 16)}

La page 16, barrée de deux traits de crayon, traite d'autre chose, dont le
dossier ne donne pas le contexte : une donnée $(\mathcal{M}, T)$ — la page
ajoute « i.e. $G$ » — munie de structures adéliques
$(T_{\mathcal{A}'}, u_{\mathcal{A}'}, t_{\infty}, t_{p})$, où $\mathcal{A}$
désigne les adèles de $\mathbb{Q}$.\footnote{Peut-être des réalisations d'un
motif sur $\mathbb{Q}$, comme aux pages 70--72 ; c'est une conjecture de
lecture, que la page ne confirme pas.} Deux énoncés :

(3) \emph{Indétermination.} La catégorie de ces données, pour
$(\mathcal{M}, T)$ fixé, est un torseur sous la catégorie des
$\mu_{2}$-torseurs sur $\mathcal{A}$ ; ses classes d'isomorphisme forment un
torseur sous $H^{1}(\mathcal{A}, \mu_{2}) = \mathcal{A}^{*}/\mathcal{A}^{*2}$,
et les automorphismes d'un objet forment $H^{0}(\mathcal{A}, \mu_{2}) =
\prod_{v} \{\pm 1\}$.\footnote{La page écrit $H^{1}(\mathcal{A}, \mu_{2})
\simeq \prod_{p} \mathbb{Q}_{p}^{*}/\mathbb{Q}_{p}^{*2}$ et
$H^{0} = \prod_{p} \mathbb{Z}/2\mathbb{Z}$. Le groupe exact est le produit
restreint des $\mathbb{Q}_{v}^{*}/\mathbb{Q}_{v}^{*2}$ relativement aux
$\mathbb{Z}_{v}^{*}/\mathbb{Z}_{v}^{*2}$, et les produits portent sur toutes
les places $v$, y compris la place réelle.}

(4) \emph{Rigidité.} Les équivalences entre deux telles données forment une
catégorie vide ou un torseur sous la catégorie des $\mu_{2}$-torseurs sur
$\mathbb{Q}$ trivialisés sur $\mathcal{A}$. Les invariants de cette dernière
sont les groupes de cohomologie relative
\[
H^{1}(\mathbb{Q} \bmod \mathcal{A}, \mu_{2}) \simeq
\Bigl(\prod_{v} \mathbb{Z}/2\mathbb{Z}\Bigr)\big/(\pm 1), \qquad
H^{0}(\mathbb{Q} \bmod \mathcal{A}, \mu_{2}) = 0 .
\]
Les deux valeurs sont justes : elles découlent de la suite exacte
$\mu_{2}(\mathbb{Q}) \to \mu_{2}(\mathcal{A}) \to H^{1}(\mathbb{Q} \bmod
\mathcal{A}) \to \mathbb{Q}^{*}/\mathbb{Q}^{*2} \to
\mathcal{A}^{*}/\mathcal{A}^{*2}$, dont la dernière flèche est injective — un
rationnel partout localement carré est un carré. C'est aussi ce qui donne la
conclusion de la page : l'image de $H^{1}(\mathbb{Q} \bmod \mathcal{A})$ dans
$H^{1}(\mathbb{Q}, \mu_{2})$ est nulle, donc l'équivalence rationnelle entre
deux données est déterminée à isomorphisme unique près, et l'on peut les
identifier sans ambiguïté.\footnote{La lecture « $\mathbb{Q} \bmod
\mathcal{A}$ » (cohomologie de $\mathbb{Q}$ relative aux adèles) est douteuse
sur la page ; elle est confortée par la notation
« $S_{\mathrm{zar}} \bmod S_{\mathrm{cris}}$ » de la page 19, qui est du même
type. Le terme médian de la formule, illisible, est omis.}

\pagerange{19}{21}
\section*{La descente comme problème de déformation (pages 19 et 21)}

\subsection*{Où est l'arbitraire (page 19)}

La page 19 prend le problème du point de vue de la descente. On part des deux
triangles $L \to \Delta(\mathcal{M}) \to L'$ et
$\mathcal{M}_{/F} \to \varphi^{*}\Delta(\mathcal{M}) \to \mathcal{M}_{/V}$,
et de la suite exacte longue de cohomologie, de la forme
\[
0 \to n \to D_{1} \to t^{*} \to \omega \to D_{0} \to \nu^{*} \to 0,
\]
où $D_{1}$, $D_{0}$ sont les cohomologies de $\Delta$. On la coupe en
\[
0 \to n \to D_{1} \to \tilde D_{1} \to 0, \qquad
0 \to \tilde D_{0} \to D_{0} \to \nu^{*} \to 0, \qquad
0 \to \tilde D_{1} \to t^{*} \to \omega \to \tilde D_{0} \to 0,
\]
les deux premières « déterminées de façon unique par $\mathrm{I}(\mathcal{M})$ ».
Ce qu'il faut descendre se réduit alors à un isomorphisme entre le complexe
$[t^{*} \to \omega]$ et un complexe $\tilde\Delta(\mathcal{M})$ d'amplitude 1
de cohomologies $\tilde D_{1}$, $\tilde D_{0}$, isomorphisme dont on connaît
l'image par $\varphi^{*}$.\footnote{La seconde suite exacte de la page,
$0 \to {}_{F}\tilde{\mathcal{M}} \to \cdots \to \tilde{\mathcal{M}}_{V} \to 0$,
est lue par morceaux, et ses indices sont incertains ; on ne la reproduit
pas. Deux tentatives, « (a) descendre $\mathcal{M}_{/F}$ et
$\mathcal{M}_{/p}$ » et « (b) descendre le triangle tout entier », sont
biffées.} Et la page situe l'arbitraire :
\[
\text{automorphismes} \in \operatorname{Ext}^{0}(S_{\mathrm{zar}} \bmod
S_{\mathrm{cris}};\ \tilde D_{0}, \tilde D_{1}) = 0,
\]
\[
\text{choix} \in \operatorname{Ext}^{1}(\ldots), \quad
\text{obstruction} \in \operatorname{Ext}^{2}(\ldots),
\]
où « $S_{\mathrm{zar}} \bmod S_{\mathrm{cris}}$ » désigne des $\operatorname{Ext}$
relatifs au foncteur $\varphi^{*}$ — ceux qui mesurent la différence entre
$\operatorname{Ext}_{S_{\mathrm{zar}}}(\tilde D_{0}, \tilde D_{1})$ et
$\operatorname{Ext}_{S_{\mathrm{cris}}}(\varphi^{*}\tilde D_{0},
\varphi^{*}\tilde D_{1})$. C'est le schéma habituel d'une théorie des
déformations, en degrés 0, 1 et 2, et le « $= 0$ ! » de la page dit que la
descente, quand elle existe, est rigide.\footnote{La page n'en donne pas de
définition, et l'indexation de ces $\operatorname{Ext}$ relatifs est notre
lecture de l'usage qu'elle en fait. Une grande partie de la prose de cette
moitié de page est illisible ; la phrase « l'arbitraire est dans \ldots{}
principalement, sans doute, i.e. dans $\operatorname{Ext}^{1}(\tilde D_{0},
\tilde D_{1})$ ou plus \ldots{} dans $\operatorname{Ext}^{1}(S_{\mathrm{zar}}
\bmod S_{\mathrm{cris}}; \ldots)$ » est ce qui en reste.}

En haut à droite de la moitié inférieure, un diagramme de schémas : $S$
plongé dans un $X$, les produits $X \times_{S} X$, $X \times_{S} X \times_{S}
X$ avec leurs enveloppes $\tilde X$, $\widetilde{X \times X}$,
$\widetilde{X \times X \times X}$, plates sur $S$, et une flèche « $F_{S}$ »
descendant de $\tilde X$ à $S$. C'est le calcul habituel d'un cristal par les
enveloppes à puissances divisées d'un plongement et de ses puissances
fibrées, auquel la factorisation du Frobenius ajoute la flèche vers
$S$.\footnote{La lecture du diagramme est approchée sur la page ; que les
$\tilde{\ }$ soient des enveloppes à puissances divisées est notre
interprétation.}

\subsection*{Une comparaison qui n'est pas injective (page 21)}

La page 21, au crayon, sans une phrase suivie, fait le calcul que la question
de fidélité de la page 11 appelle. Elle munit $\Delta$ de la filtration
$\operatorname{Fil}^{0} = \Delta$, $\operatorname{Fil}^{1} = \ell$,
$\operatorname{Fil}^{2} = 0$, et compare, pour $\Theta =
\mathcal{O}_{X}[1]$,
\[
\operatorname{Hom}_{D(X)}(\check\ell'[1], \Theta) \to
\operatorname{Hom}(\Delta, \Theta)
\quad\text{et}\quad
\operatorname{Hom}_{D(X_{\mathrm{cris}})}(\varphi^{*}\check\ell'[1],
\varphi^{*}\Theta).
\]
Elle écrit la suite exacte à six termes
\[
0 \to n_{G^{*}} \to \mathbb{D}^{*}(G)^{\vee} \to t_{G} \to \omega_{G^{*}}
\to \mathbb{D}^{*}(G^{*}) \to \nu_{G} \to 0,
\]
qui est la cohomologie du triangle (III) écrit pour $G^{*}$, et la suite
d'hypercohomologie correspondante. Elle conclut par deux constats :
$H^{1}(\mathbb{D}^{*}(G)^{\vee}) = 0$, et
\[
\operatorname{Ker}\bigl(H^{1}(n_{G^{*}}) \to H^{1}_{\mathrm{cris}}(n_{G^{*}})\bigr)
\ne 0 .
\]
Autrement dit : la cohomologie de Zariski ne s'injecte pas dans la cohomologie
cristalline de l'image par $\varphi^{*}$, et c'est un obstacle à la fidélité
espérée page 11.\footnote{La dernière lettre de la suite à six termes,
« $\nu_{G}$ », est une lecture incertaine, de même que les accents et
exposants de la ligne $0 \to \ell_{\bullet} \to \check\Delta^{*} \to
\check\ell_{\bullet}[1] \to 0$. La page ne dit pas sous quelles hypothèses le
noyau est non nul ; le rapport avec la question de la page 11 est notre
lecture.}

\pagerange{24}{30}
\section*{Groupes de Barsotti--Tate ordinaires et $F$-cristaux à connexion (pages 24 à 30)}

\subsection*{La catégorie $\mathbf{C}$ (page 24)}

Soit $S_{0}$ un schéma lisse sur un corps parfait $k$, $S$ un schéma formel
formellement lisse sur $W$ qui le relève, et $F_{S}$ un endomorphisme de $S$
relevant le Frobenius absolu de $S_{0}$ et compatible à $\sigma$ sur $W$.
Soit $\mathrm{BTO}(S_{0})$ la catégorie des groupes de Barsotti--Tate
ordinaires sur $S_{0}$, et $\mathbf{C} = \mathbf{C}_{S, F_{S}}$ la catégorie
des sextuples $(P', Q', M, \nabla, F_{M}, \operatorname{Fil}_{0})$ :

\begin{itemize}
\item $P'$, $Q'$ des faisceaux $p$-adiques localement constants (tordus),
libres de type fini ;
\item $M$ une extension $0 \to Q' \otimes_{\mathbb{Z}_{p}} \mathcal{O}_{S}
\to M \to P' \otimes_{\mathbb{Z}_{p}} \mathcal{O}_{S} \to 0$ ;
\item $\nabla$ une connexion intégrable sur $M$, compatible à la structure
d'extension, les termes extrêmes ayant leur connexion canonique ;
\item $F_{M}\colon F_{S}^{*}M \to M$ un morphisme d'extensions horizontal,
induisant l'isomorphisme canonique sur $Q'$ et $p$ fois l'isomorphisme
canonique sur $P'$ ;
\item $\operatorname{Fil}_{0}$ un scindage de l'extension réduite $M_{0}$,
compatible à $F_{M_{0}}$.
\end{itemize}

Sur $Q'$ le Frobenius est un isomorphisme (pente 0), sur $P'$ il est $p$ fois
un isomorphisme (pente 1) : c'est ce qu'« ordinaire » veut dire. On a un
foncteur canonique, le module de Dieudonné,
$\mathbb{D}\colon \mathrm{BTO}(S_{0})^{\circ} \to \mathbf{C}$, et la page
pose la \emph{Question} : est-ce une équivalence de catégories ?\footnote{La
question n'est pas résolue dans le dossier. La pleine fidélité du foncteur de
Dieudonné cristallin sur une base lisse sur un corps parfait a été établie
par de Jong (1995) ; on ne prétend pas qu'elle réponde à la question telle
qu'elle est posée ici, avec cette catégorie $\mathbf{C}$.}

\subsection*{En coordonnées (pages 25 et 26)}

Choisissons un scindage de l'extension,
$s\colon M \xrightarrow{\sim} P'\otimes\mathcal{O}_{S} \oplus
Q'\otimes\mathcal{O}_{S}$, et notons $T = P'^{\vee} \otimes_{\mathbb{Z}_{p}} Q'
= \mathcal{H}om(P', Q')$. Un autre scindage s'écrit $s' = u\,s$ avec
\[
u(x, y) = (x,\ y + Lx), \qquad L \in \Gamma(T \otimes_{\mathbb{Z}_{p}}
\mathcal{O}_{S}).
\]
Le Frobenius, lu dans le scindage $s$, est
\[
F_{M,s}(x, y) = (px,\ y + \Phi x), \qquad \Phi = \Phi_{M,s} \in
\Gamma(T \otimes_{\mathbb{Z}_{p}} \mathcal{O}_{S}),
\]
et la connexion est $\nabla_{M,s}(x, y) = (dx,\ dy + \omega x)$, avec
$\omega = \omega_{M,s} \in \Gamma(\hat\Omega^{1}_{S/W} \otimes T)$.\footnote{La
page écrit la matrice de $F_{M,s}$ avec $\mathrm{id}$ en haut et
$p\,\mathrm{id}$ en bas, et la formule $(x,y) \mapsto (px, y + \Phi x)$ ; les
deux ne s'accordent pas. On suit la formule, qui est celle qu'exige la
définition de $\mathbf{C}$ ($p$ sur $P'$, l'identité sur $Q'$). Le $T$ est
écrit sur un $\check P$.} De $F_{M,s'} F_{S}^{*}(u) = u F_{M,s}$ et de
$\nabla_{M,s'} = u\nabla_{M,s}u^{-1}$ on tire les lois de changement de
scindage
\[
\Phi' = \Phi + \bigl(pL - F_{S}^{*}(L)\bigr), \qquad \omega' = \omega - dL .
\]
Et le triple de conditions qui exprime qu'on a bien un objet de $\mathbf{C}$ :
\[
d\omega = 0, \qquad F_{S}^{*}(\omega) - p\,\omega = d\Phi, \qquad
\Phi_{0} = 0,
\]
la première exprimant que $\nabla$ est intégrable (le terme
$\omega \wedge \omega$ est nul, $\omega$ étant strictement triangulaire), la
deuxième que $F_{M}$ est horizontal, la troisième que $F_{M_{0}}$ respecte
$\operatorname{Fil}_{0}$.\footnote{On a refait les trois calculs ; ils sont
justes. La troisième condition suppose le scindage $s$ choisi de façon à
relever $\operatorname{Fil}_{0}$ sur $S_{0}$, ce que dit une note de marge de
la page 25 : on peut prendre un tel scindage, « ce qui conduira à des $L$
\ldots{} \emph{nuls} sur $S_{0}$ et des $\Phi$ itou ».} La classification de
$\mathbf{C}$ est donc celle des paires $(\Phi, \omega)$ soumises à ces
conditions, modulo l'action des $L$ nuls sur $S_{0}$.

\subsection*{Le cas d'un module libre quelconque (pages 28 et 30)}

La page 28 refait le calcul pour un $F$-cristal $M = E \otimes
\mathcal{O}_{S}$, $E$ libre de type fini, de matrice de connexion
$A \in \Gamma(\hat\Omega^{1}_{S/W} \otimes \operatorname{End} E)$ et de
Frobenius $\Phi\colon F_{S}^{*}(E_{S}) \simeq E_{S} \to E_{S}$, qui devient
un isomorphisme après $\otimes \mathbb{Q}_{p}$. Avec $A^{\sigma} =
F_{S}^{*}(A)$ :
\[
d\Phi = \Phi A^{\sigma} - A\,\Phi, \qquad dA + A \wedge A = 0, \qquad
\exists\,\Psi,\ i : \ \Psi\Phi = p^{i}\,\mathrm{id},
\]
pour l'horizontalité de $F$, l'intégrabilité, et le fait que $\Phi$ est une
isogénie.\footnote{La page écrit la condition d'intégrabilité
$dA + [A, A] = 0$. Avec la convention usuelle $[A, A] = A\wedge A + A \wedge
A$ pour les formes à valeurs matricielles, la condition juste est
$dA + A \wedge A = 0$, soit $dA + \tfrac12[A, A] = 0$. En tête de page, au
crayon, un calcul sans rapport apparent, $(i+p)(p+1) - (i+p-1)p = 2p + i$,
qui est juste, et un encadré « $v \circ \omega = du + \omega' \circ u$ » où
l'on attend $u$ au lieu de $v$ : c'est la condition de la page 30.} La page
30 donne les morphismes : $u\colon E_{S} \to E'_{S}$ est un morphisme si
\[
du = u A - A' u, \qquad F_{M'}\, u^{\sigma} = u\, F_{M} .
\]

\emph{Exemple : les modules inversibles.} Pour $M$ de rang 1, $\Phi =
a \in \Gamma(\mathcal{O}_{S})$ et $A = \omega \in
\Gamma(\hat\Omega^{1}_{S/W})$ ; les conditions deviennent
\[
da = a(\omega^{\sigma} - \omega), \qquad d\omega = 0,
\qquad ab = p^{i},
\]
et deux tels objets $(a, \omega)$, $(a', \omega')$ sont
isomorphes s'il existe $u \in \Gamma(\mathcal{O}_{S}^{*})$ avec
\[
\omega' - \omega = -\frac{du}{u}, \qquad \frac{a'}{a} =
\frac{u}{u^{\sigma}} .
\]
D'où la classification des données de rang 1 sur un module trivial par la
cohomologie du complexe
\[
\Gamma(\mathcal{O}_{S}^{*}) \longrightarrow \Gamma(\hat\Omega^{1}_{S/W})
\times \Gamma(\mathcal{O}_{S}[1/p])^{*} \longrightarrow
\Gamma(\hat\Omega^{2}_{S/W}) \times \Gamma(\hat\Omega^{1}_{S/W})[1/p],
\]
\[
u \longmapsto \Bigl(-\frac{du}{u},\ \frac{u}{u^{\sigma}}\Bigr), \qquad
(\omega, a) \longmapsto \Bigl(d\omega,\
\frac{da}{a} - (\omega^{\sigma} - \omega)\Bigr),
\]
le $H^{0}$ donnant les automorphismes et le $H^{1}$ les classes
d'isomorphisme.\footnote{On renomme $a$ et $b$ ce que la page appelle
$\varphi$ et $\psi$, la première lettre étant réservée. La page écrit la première flèche
$u \mapsto (du/u,\ u/u^{\sigma})$ ; avec ce signe le composé des deux
flèches n'est pas nul (il vaut $2(du/u - F_{S}^{*}(du/u))$ dans la seconde
composante). Le signe $-du/u$, qui est celui de sa propre loi de changement
$\omega' - \omega = -du/u$, rend le composé nul. Le « ! » après
$\Gamma(\mathcal{O}_{S})$ et la glose « éléments inversibles dans
$A[1/p]$ » sont sur la page ; l'exposant de $p$ dans $\varphi\psi = p^{i}$
est une lecture douteuse.}

\pagerange{29}{29}
\subsection*{Une page de topologie différentielle (page 29)}

La page 29, au stylo à bille, est d'un autre sujet : des familles d'ouverts
$(\sigma_{i}, (U_{i}, V_{i}))$ de classe $C^{r}$ adaptées à une
sous-variété, $\sigma(V) = \sigma(U) \cap \mathbb{R}^{n}$, avec une
hiérarchie de familles « canoniques », « régulières », « naturelles »,
« normales », et les conditions de compatibilité de deux cartes. Elle est
sans rapport avec les cristaux et proche, par le sujet, du cours polycopié
sur les variétés $C^{r}$ dont des feuillets servent de versos dans ce dossier.

\pagerange{31}{31}
\subsection*{Un feuillet de tapuscrit (page 31)}

Un feuillet de tapuscrit (numéroté 34) porte la démonstration d'un énoncé à
trois volets sur la réduction semi-stable des variétés abéliennes —
critère galoisien 3.5, réduction au cas strictement local, trivialité de la
représentation de Galois sur les trois quotients de la filtration de
monodromie, calculée par sa trace — puis la Remarque 5.13.1 sur les groupes de
Barsotti--Tate à bonne réduction virtuelle. Les interventions à la main sont
des lettres tracées dans les blancs, « (pro-$\ell$-) » ajouté devant
« groupe de Barsotti--Tate », et à la fin la parenthèse « (C'est vrai pour
$\ell \ne p$.) ».\footnote{Le « critère galoisien de réduction semi-stable
3.5 » évoque l'exposé IX de SGA 7 ; non vérifié. Ce feuillet n'appartient pas
au tapuscrit des pages 6 à 14.}

\pagerange{33}{41}
\section*{Les groupes de Barsotti--Tate à multiplication par $A$ (pages 33 à 41)}

\subsection*{Le module de Dieudonné (pages 33 et 35)}

On veut déterminer les groupes de Barsotti--Tate $G_{0}$ sur $k =
\mathbb{F}_{q}$, de hauteur $h$, sur lesquels $A$ opère. Par la théorie de
Dieudonné, ce sont les $W$-modules libres $M$ de rang $h$ munis d'un
endomorphisme $\sigma$-linéaire $F_{M}$ et d'une action de $A$ par
endomorphismes commutant à $F_{M}$, avec
\[
(*') \qquad \exists\, V_{M} : \ F_{M}V_{M} = V_{M}F_{M} = p .
\]
Le cas particulier qui intéresse la page est celui où l'action de $\pi$ est
le Frobenius relatif à $\mathbb{F}_{q}$ :
\[
(*) \qquad F_{M}^{f} = \pi_{M} .
\]
Les deux actions font de $M$ un module sur $A_{W} = A \otimes_{\mathbb{Z}_{p}}
W$, et $F_{M}$ est $\sigma_{A}$-linéaire pour $\sigma_{A} = \mathrm{id}
\otimes \sigma$. Comme $k$ est le corps résiduel de $K$, $A$ contient $W$
et
\[
A \otimes_{\mathbb{Z}_{p}} W \xrightarrow{\ \sim\ } A^{\Sigma}, \qquad
\Sigma = \operatorname{Gal}(k/\mathbb{F}_{p}) \simeq \mathbb{Z}/f\mathbb{Z},
\]
$\sigma_{A}$ permutant circulairement les facteurs. Un $A_{W}$-module $M$ est
donc une famille $(M_{\nu})_{\nu \in \mathbb{Z}/f}$ de $A$-modules, et
$F_{M}$ une famille d'applications $A$-linéaires
\[
M_{0} \xrightarrow{F_{0}} M_{1} \xrightarrow{F_{1}} \cdots \to M_{f-1}
\xrightarrow{F_{f-1}} M_{0} .
\]
Par $(*')$, $F_{M}$ est injectif, donc chaque $F_{\nu}$ l'est, et les
$M_{\nu}$ ont tous le même rang. Comme les rangs des $M_{\nu}$ \emph{sur
$W$} ont pour somme $h = ef$, chacun est de rang $e$ sur $W$, c'est-à-dire de
rang 1 sur $A$ : les $M_{\nu}$ sont inversibles.\footnote{La page invoque la
condition $(*)$ pour l'injectivité, et la marge note que $(*')$ seule
suffirait ; c'est le cas. Elle dit « la somme des rangs des $M_{\nu}$ est
$h$ » : c'est vrai des rangs sur $W$ ; la somme des rangs sur $A$ est $f$.
Les composantes du Frobenius sont notées $\varphi_{\nu}$ sur la page.}

On identifie alors tous les $M_{\nu}$ à $\Omega = M_{0}$ par les
isomorphismes $\pi^{-m_{\nu}} F_{\nu-1}\cdots F_{0}\colon M_{0} \to M_{\nu}$,
où $0 = m_{0} \le m_{1} \le \cdots \le m_{f-1}$ sont les entiers uniquement
déterminés qui rendent ces flèches surjectives. On obtient un isomorphisme
canonique
\[
M \simeq \Omega \otimes_{\mathbb{Z}_{p}} W = \Omega \otimes_{A}
(A \otimes_{\mathbb{Z}_{p}} W).
\]

\subsection*{La dimension (page 37)}

Dans ces identifications, $F_{\nu} = \pi^{r_{\nu}}\,\mathrm{id}_{\Omega}$
pour $0 \le \nu \le f-2$, avec $r_{\nu} = m_{\nu+1} - m_{\nu} \ge 0$, et
$F_{f-1} = \pi^{r_{f-1}}u$ avec $u \in A^{*}$. Donc
\[
F_{M}^{f} = \varpi_{M}, \qquad \varpi = \pi^{r_{0} + \cdots + r_{f-1}}\,u \in A,
\]
la condition $(*')$ s'écrit $r_{\nu} \le e$ pour tout $\nu$ (chaque
$F_{\nu}$ divise $p$, de valuation $e$), et
\[
M/F_{M}M \simeq \bigoplus_{\nu} A/\pi^{r_{\nu}}A, \qquad
\dim G_{0} = r_{0} + \cdots + r_{f-1} = v_{\pi}(\varpi) \ \le\ fe = h .
\]
C'est la formule encadrée de la page.\footnote{Que $\dim G_{0}$ soit la longueur de $M/F_{M}M$ (et non de
$M/V_{M}M$) dépend de la convention de Dieudonné, ici contravariante ; avec
l'autre convention on obtient la dimension du dual. Le module est un
$W$-module de longueur $\sum r_{\nu}$ parce que $A$ et $W$ ont même corps
résiduel. La page a la même formule, avec un premier membre biffé
illisible.}

\subsection*{La proposition (page 39)}

\emph{Proposition.} Soit $G_{0}$ un groupe de Barsotti--Tate de hauteur $h$
sur $k$ sur lequel $A$ opère, et $M$ son module de Dieudonné. Alors
$F_{G_{0}}^{f}$ est de la forme $\varpi_{G_{0}}$ pour un $\varpi \in A$, et
$\dim G_{0} = \operatorname{long}_{A} A/\varpi A = v(\varpi)$. Les conditions
suivantes sont équivalentes :
\begin{enumerate}
\item[(i)] $\dim G_{0} = 1$ ;
\item[(ii)] $F_{G_{0}}^{f} = \varpi_{G_{0}}$ avec $\varpi$ une uniformisante
de $A$ — uniquement déterminée ;
\item[(iii)] $M$ est donné par un triple $(\Omega, \nu, \varpi)$, $\Omega$
un $A$-module inversible, $0 \le \nu \le f-1$, $\varpi$ une uniformisante :
\[
M = \Omega \otimes_{\mathbb{Z}_{p}} W, \qquad
F_{M} = \mathrm{id}_{\Omega} \otimes \sigma + \mathrm{pr}_{\nu}\bigl[(\varpi - 1)
\otimes \sigma\bigr],
\]
où $\mathrm{pr}_{\nu}$ est le projecteur sur le $\nu$-ième facteur de
$A \otimes_{\mathbb{Z}_{p}} W \simeq A^{\Sigma}$.
\end{enumerate}

La démonstration est le calcul précédent : $v(\varpi) = 1$ force un seul
$r_{\nu_{0}}$ égal à 1 et les autres nuls, et en modifiant les
identifications $M_{0} \simeq M_{\nu}$ par l'unité $u$ pour $\nu > \nu_{0}$
on fait passer $u$ de $F_{f-1}$ à $F_{\nu_{0}}$, qui devient $\varpi =
\pi u$.\footnote{La page note, en petit, « ($u = 1$ \ldots{} $\varpi = \pi$) »
et « quitte à changer $\pi$ ! ». Le passage de $u$ d'une composante à
l'autre est notre explicitation. Une note de marge, en grande partie
illisible, affirme que tous ces groupes sont isomorphes par des
transformations faisant intervenir $f$ et $\mathbb{F}_{q} = \mathbb{F}_{p^{f}}$ ;
on ne peut pas la reconstituer. Le paramètre $\nu$ mesure, d'après la page
43, le décalage entre deux structures de $k$-espace sur $\omega_{G_{0}}$, et
n'est donc pas un invariant de $G_{0}$ comme groupe à opérations de $A$
qu'on puisse faire disparaître sans tordre l'action.}

\subsection*{Les homomorphismes (pages 39 et 41)}

Soient $G_{0}$, $G'_{0}$ comme plus haut, donnés par $(\Omega, (r_{\nu}), u)$
et $(\Omega', (r'_{\nu}), u')$. Un homomorphisme $A$-linéaire $G_{0} \to
G'_{0}$ est un morphisme $M' \to M$ de $A_{W}$-modules commutant au
Frobenius. On en trouve un non nul — c'est alors une isogénie — si et
seulement si $\varpi = \varpi'$, c'est-à-dire $F_{G_{0}}^{f} =
F_{G'_{0}}^{f}$ dans $A$. Dans ce cas, le $A$-module de ces homomorphismes
est isomorphe au sous-module $H \subset \operatorname{Hom}(\Omega', \Omega)$
des $v_{0}$ tels que
\[
\operatorname{val}(v_{0}) + \sum_{i=0}^{\nu} r_{i} - \sum_{i=0}^{\nu} r'_{i}
\ \ge\ 0 \qquad (0 \le \nu \le f-2) :
\]
la commutation au Frobenius impose $v_{\nu+1} = \pi^{\sum_{0}^{\nu}(r_{i} -
r'_{i})}v_{0}$ sur la composante $\nu+1$, et la condition de la composante
0, où le cycle se referme, est automatique dès que $\varpi = \varpi'$.

En dimension 1, avec $r_{i} = \delta_{i\nu}$ et $r'_{i} = \delta_{i\nu'}$,
cela donne :
\[
\operatorname{Hom}_{A}(G_{0}, G'_{0}) \simeq
\begin{cases}
\operatorname{Hom}(\Omega', \Omega) & \text{si } \nu \le \nu', \\
\pi\operatorname{Hom}(\Omega', \Omega) & \text{si } \nu > \nu',
\end{cases}
\]
et $G_{0}$, $G'_{0}$ sont isogènes si et seulement si leurs uniformisantes
coïncident.\footnote{Le signe entre $\nu$ et $\nu'$ dans la première ligne
est couvert d'une tache d'encre ; le calcul ci-dessus donne $\le$. Les primes
sur les $r_{i}$ des deux sommes de la page 39 ne se distinguent pas
nettement ; leur place est fixée ici par le calcul.}

\subsection*{À isogénie près (page 41)}

Sans supposer $\dim G_{0} = 1$, à isogénie près les $M_{\nu}$ deviennent des
$K$-espaces de dimension 1, tous identifiés par les $F_{\nu}$, et
\[
M \otimes \mathbb{Q} \simeq \Omega \otimes_{\mathbb{Q}_{p}} K_{0} =
\Omega \otimes_{K} (K \otimes_{\mathbb{Q}_{p}} K_{0}), \qquad
F_{M}^{f} = \varpi \otimes \mathrm{id}_{K_{0}},
\]
avec $\Omega$ un $K$-espace de dimension 1 et $\varpi \in K^{*}$. La
condition d'« effectivité » — l'existence d'un réseau stable par $F$ et par
$V = pF^{-1}$ — entraîne que $\varpi$ et $V^{f} = p^{f}\varpi^{-1}$
préservent un réseau, c'est-à-dire $0 \le v_{\pi}(\varpi) \le fe$ ; et
réciproquement cette condition suffit, puisqu'on peut alors répartir
$v_{\pi}(\varpi)$ en $f$ entiers $r_{\nu} \in [0, e]$.\footnote{La page écrit
la condition comme l'existence de réseaux invariants par $\varpi$ et par
$p^{f}\varpi^{-1}$ ; l'exposant $f$ est une lecture douteuse, confirmée par
$V^{f} = p^{f}F^{-f}$. La suffisance, qui demande un réseau stable par $F$ et
$V$ eux-mêmes et non par leurs puissances $f$-ièmes, est notre complément.}

\pagerange{43}{49}
\section*{Le relèvement de Lubin--Tate (pages 43 à 49)}

\subsection*{L'espace cotangent et le groupe $G_{0}(A, \pi)$ (page 43)}

Sous les conditions de la proposition, $\omega_{G_{0}}$ est la seule
composante non nulle de $M/F_{M}M$, isomorphe à $\Omega/\pi\Omega$. Il porte
deux structures de $k$-espace : celle qui vient de $A/\pi A = k$ par
l'action de $A$, et celle qui vient de $W/pW = k$. Elles diffèrent d'une
puissance de $\sigma$ déterminée par $\nu$, et le cas le plus simple, $\nu =
0$, est celui où elles coïncident.\footnote{Trois lignes surchargées de la
page, en partie biffées, décrivent la filtration canonique de
$M \otimes_{W} k = \sum M_{\nu} \otimes_{A} k$ ; seuls des fragments se
lisent. L'exposant de $\pi$ dans « $M_{\nu}/\pi^{r}M_{\nu}$ » et la lettre
devant le dernier $\Omega$ sont douteux. Quelle composante porte l'indice 0,
et donc quelle valeur de $\nu$ fait coïncider les deux structures, dépend de
conventions d'indexation que la page ne fixe pas ; on suit son énoncé.} Avec
de plus $\Omega = A$, on obtient le groupe
\[
G_{0} = G_{0}(A, \pi), \qquad M = A \otimes_{\mathbb{Z}_{p}} W, \qquad
F_{M} = \mathrm{id}_{A} \otimes \sigma + \mathrm{pr}_{0}\bigl((\pi - 1)
\otimes \sigma\bigr),
\]
où $\mathrm{pr}_{0}$ est le projecteur sur le facteur correspondant au
$A$-homomorphisme $A \otimes_{\mathbb{Z}_{p}} W \to A$ qui est l'inclusion
sur $W$ ; le cas d'un $\Omega$ quelconque s'en déduit par le produit
tensoriel $\Omega \otimes_{A} G_{0}(A)$.

\subsection*{Le relèvement par la théorie des déformations (pages 43 à 47)}

On veut relever $G_{0}$ en un groupe de Barsotti--Tate $G$ sur $A$, muni
d'une action de $A$, de sorte que l'action induite de $A$ sur le $A$-module
inversible $t_{G}$ soit l'action évidente. D'après Lubin--Tate, un tel
relèvement existe et est unique à isomorphisme unique près.\footnote{Lubin
et Tate, « Formal complex multiplication in local fields » (1965). Le groupe
cherché est un $A$-module formel au sens strict.} La page en redonne une
preuve, sous une hypothèse : que l'idéal $\pi A$ admette des puissances
divisées, ce qui a lieu si et seulement si $e \le p-1$.\footnote{La page dit
« $\pi A$ stable par puissances divisées » ; la condition $e \le p - 1$ est la
traduction usuelle. La théorie des déformations des groupes de
Barsotti--Tate par le cristal de Dieudonné et la filtration de Hodge est
celle que Grothendieck annonce à Nice en 1970 et que Messing rédige (1972) ;
on la cite sous son nom d'aujourd'hui, Grothendieck--Messing.}

Sous cette hypothèse, relever $G_{0}$ à isomorphisme près revient à relever
la filtration de Hodge de $M \otimes_{W} k$ en une filtration de
$M_{A} = \Omega \otimes_{\mathbb{Z}_{p}} A$. La compatibilité à l'action de
$A$ veut dire que cette filtration est faite de
$A \otimes_{\mathbb{Z}_{p}} A$-modules : on cherche
\[
0 \to \omega_{G} \to \Omega \otimes_{\mathbb{Z}_{p}} A \to t_{G^{*}} \to 0,
\]
ou, par transposition,
\[
(**) \quad 0 \to \omega_{G^{*}} \to \check\Omega \otimes_{\mathbb{Z}_{p}} A
\to t_{G} \to 0 .
\]
Que $A$ opère sur $t_{G}$ par sa structure évidente veut dire que les deux
actions de $A$ sur $t_{G}$, par les deux facteurs de $A \otimes A$,
coïncident, c'est-à-dire que $A \otimes A$ opère sur $t_{G}$ à travers la
multiplication $p_{0}\colon A \otimes_{\mathbb{Z}_{p}} A \to A$. Comme $t_{G}$
est de rang 1 sur $A$, et $\check\Omega \otimes_{\mathbb{Z}_{p}} A$ de rang 1
sur $A \otimes A$, et que $p_{0}$ est surjectif, $t_{G}$ est le changement de
base de $\check\Omega \otimes_{\mathbb{Z}_{p}} A$ par $p_{0}$ ; et comme
$\check\Omega \otimes_{\mathbb{Z}_{p}} A$ se déduit de $\check\Omega$ par
$\lambda \mapsto \lambda \otimes 1$, dont le composé avec $p_{0}$ est
l'identité,
\[
t_{G} \simeq \check\Omega, \qquad \omega_{G} \simeq \Omega .
\]
Les suites $(**)$ et sa duale sont ainsi canoniques et prolongent la
filtration de $M_{A} \otimes_{A} k$ : c'est l'existence et l'unicité de
Lubin--Tate, sous une forme plus faible quant à l'unicité, et sous
l'hypothèse supplémentaire des puissances divisées.\footnote{La flèche
$p_{0}$ de $(**)$ est écrite sur une lettre surchargée.}

\subsection*{Sans puissances divisées : à isogénie près (pages 47 et 49)}

Si l'on abandonne l'hypothèse et qu'on utilise la classification des
déformations à isogénie près, on trouve de même des suites exactes
canoniques sur $K$,
\[
0 \to \omega_{G} \otimes_{A} K \to M_{K} \to t_{G^{*}} \otimes_{A} K \to 0,
\qquad
0 \to \omega_{G^{*}} \otimes_{A} K \to \check M_{K} \to t_{G} \otimes_{A} K
\to 0,
\]
avec $M_{K} = \Omega \otimes_{\mathbb{Z}_{p}} K$, et $t_{G} \otimes_{A} K
\simeq \check\Omega \otimes_{A} K$ : l'unicité à isogénie près, mais ni
l'existence ni l'unicité à isomorphisme près, pour lesquelles il faut
revenir à Lubin--Tate.\footnote{L'indice du produit tensoriel médian de la
formule de la page 47 est une lecture douteuse.}

Les deux réseaux $t_{G}$ et $\check\Omega$ d'une même droite sur $K$
diffèrent alors d'une puissance de $\pi$ :
\[
t_{G} = \pi^{-\alpha}\check\Omega, \qquad \text{donc} \qquad
\omega_{G} = \pi^{\alpha}\Omega, \qquad \alpha \in \mathbb{Z},
\]
$\alpha$ ne dépendant que de $A$ et de $\pi$.\footnote{La page écrit
« $\Omega = \pi^{\alpha}\omega_{G}$ ». Le dual de $\pi^{-\alpha}\check\Omega$
est $\pi^{\alpha}\Omega$, de sorte que $\Omega = \pi^{-\alpha}\omega_{G}$ ;
c'est le signe qu'on corrige. La parenthèse qui devait dire de quoi dépend
$\alpha$ est illisible.} De même, le module de Dieudonné de $G$ et $M$ sont
deux réseaux du même $K$-espace $\Omega \otimes_{\mathbb{Z}_{p}} K$, et la
page pose la question de savoir s'ils coïncident, comme c'est le cas en
présence de puissances divisées — ce qui forcerait $\alpha = 0$. « Cela ne
dépend pas du choix de $\pi$, mais seulement de $A$ \ldots{} » ; la question
reste ouverte.

\pagerange{51}{55}
\section*{Le théorème de Tate pour le groupe de Lubin--Tate (pages 51 à 55)}

Soit $G$ le groupe de Lubin--Tate, et $K_{\pi}$ l'extension abélienne de $K$
engendrée par ses points de torsion, de groupe $A^{*}$ via le caractère de
Lubin--Tate $\chi_{\pi}\colon \Gamma_{K} \to A^{*}$. Le module de Tate
$T_{p} = T_{p}(G)(K_{\pi})$ est un $A$-module inversible dont les bases sont
permutées de façon simplement transitive par $\operatorname{Gal}(K_{\pi}/K)$
: il y a donc un choix privilégié, à isomorphisme unique près, de $K_{\pi}$
muni d'un générateur $e$ de $T_{p}$, et $\operatorname{Spec} K_{\pi}$
représente le faisceau des générateurs de $T_{p}(G)$. Ce choix fait,
$T_{p} \simeq A$, $\lambda \mapsto \lambda e$, et l'on prend pour $\bar K$
une extension de $K_{\pi}$.\footnote{La phrase « Nous supposons fixé
$K_{\pi}$ \ldots{} » est de lecture difficile. L'énoncé sur le faisceau des
générateurs est à entendre pour le pro-objet $\varprojlim G[\pi^{n}]^{*}$.}

On regarde sur $T_{p} \otimes_{\mathbb{Z}_{p}} \mathbb{C}_{p}$ trois actions
qui commutent : celle de $A$ par les opérations sur $G$, celle de $A$ par la
structure de $\mathbb{C}_{p}$-espace (via $K \subset \mathbb{C}_{p}$), et
celle de $\Gamma_{K}$. On suppose $K/\mathbb{Q}_{p}$ galoisienne. Comme
$A \otimes_{\mathbb{Z}_{p}} A$-module, on a canoniquement
\[
(*) \qquad T_{p} \otimes_{\mathbb{Z}_{p}} \mathbb{C}_{p} \simeq
A \otimes_{\mathbb{Z}_{p}} \mathbb{C}_{p} \simeq
\prod_{\sigma} \mathbb{C}_{p,\sigma},
\qquad \sigma \in \operatorname{Hom}_{\mathbb{Q}_{p}}(K, \mathbb{C}_{p})
\simeq \operatorname{Gal}(K/\mathbb{Q}_{p}),
\]
où $\mathbb{C}_{p,\sigma} = \mathbb{C}_{p}$ sur lequel $A$ opère via
$\sigma$. Sur le facteur $\mathbb{C}_{p,\sigma}$, $g \in \Gamma_{K}$ opère
par
\[
x \longmapsto \sigma(\chi_{\pi}(g))\, g(x) .
\]
Cette action est semi-linéaire.\footnote{La page écrit $\rho_{\sigma}(g)x = \sigma(\chi_{\pi}(g))x$, sans le
$g(x)$ : l'action est semi-linéaire, et c'est bien sous cette forme que la
page 55 l'utilise. L'identification
$\operatorname{Hom}_{\mathbb{Q}_{p}}(K, \mathbb{C}_{p}) \simeq
\operatorname{Gal}(K/\mathbb{Q}_{p})$ dépend du plongement fixé, qui devient
$\sigma = 1$.}

Le théorème de Tate donne de son côté une décomposition
\[
(**) \qquad T_{p} \otimes \mathbb{C}_{p} \simeq
\bigl(t_{G} \otimes_{A} \mathbb{C}_{p}(1)\bigr) \oplus
\bigl(\omega_{G^{*}} \otimes_{A} \mathbb{C}_{p}\bigr),
\]
fonctorielle, donc compatible aux deux actions de $A$.\footnote{Tate,
« $p$-divisible groups » (1967). La page note les deux termes
$\mathfrak{t} \otimes C(1)$ et $\check{\mathfrak{t}}^{*} \otimes C$ ; la lettre
lue $\mathfrak{t}$ est une majuscule cursive de lecture douteuse, que l'on
identifie à l'espace tangent $t_{G}$, et $\check{\mathfrak{t}}^{*}$ à
$\omega_{G^{*}} = t_{G^{*}}^{\vee}$, ce que vérifie le cas de
$\mu_{p^{\infty}}$ et de $\mathbb{Q}_{p}/\mathbb{Z}_{p}$.} Sur $t_{G}$, $A$
opère par l'action régulière : $t_{G} \otimes_{A} \mathbb{C}_{p}$ est donc le
facteur $\sigma = 1$, et l'autre terme est le produit des
$\mathbb{C}_{p,\sigma}$ pour $\sigma \ne 1$. Se donner les isomorphismes du
théorème de Tate, c'est alors se donner, pour chaque $\sigma$, un
isomorphisme $\lambda \mapsto \lambda x_{\sigma}$ ($x_{\sigma} \in
\mathbb{C}_{p}^{*}$) qui entrelace les actions de Galois, la source étant
tordue par $\sigma(\chi_{\pi})$, le but par le caractère cyclotomique
$\chi_{\mathrm{cyc}}$ si $\sigma = 1$ et non tordu sinon. Cela revient aux
conditions
\[
g(x_{\sigma}) = \sigma(\chi_{\pi}(g))\, x_{\sigma} \quad (\sigma \ne 1), \qquad
g(x_{1}) = \chi_{\mathrm{cyc}}(g)^{-1}\chi_{\pi}(g)\, x_{1}, \qquad
g \in \Gamma_{K} ,
\]
qui sont celles que la page encadre.\footnote{La page appelle $\tau$ le caractère cyclotomique (« caractère de
Tate »). On a refait le calcul de l'entrelacement ; les deux conditions sont
celles de la page.}

\emph{N.B.} A priori les $x_{\sigma}$ ne sont déterminés qu'à multiplication
près par des éléments de $K^{*}$, mais le théorème de Tate détermine tous les
autres dès qu'on a choisi l'un d'eux. Et il n'existe pas de $x'_{1} \in
\mathbb{C}_{p}^{*}$ tel que $g(x'_{1}) = \chi_{\pi}(g)\,x'_{1}$, même pour
$g$ dans l'inertie : le produit de $x'_{1}$ et des $x_{\sigma}$, $\sigma \ne
1$, serait un $x \in \mathbb{C}_{p}^{*}$ avec
\[
g(x) = N_{K/\mathbb{Q}_{p}}(\chi_{\pi}(g))\, x = \chi_{\mathrm{cyc}}(g)\,x
\qquad (g \text{ dans l'inertie}),
\]
ce qui est impossible par le théorème de Tate ($\mathbb{C}_{p}(1)$ n'a pas
d'invariant non nul sous l'inertie). En langage d'aujourd'hui : $\chi_{\pi}$
est un caractère de Hodge--Tate, de poids 1 au plongement $\sigma = 1$ et 0
aux autres, et non de poids 0.\footnote{L'égalité
$N_{K/\mathbb{Q}_{p}} \circ \chi_{\pi} = \chi_{\mathrm{cyc}}$ sur l'inertie
est la théorie du corps de classes local de Lubin--Tate ; la page l'écrit
sous l'accolade comme $N_{A^{*}/\mathbb{Z}_{p}^{*}}(\chi_{\pi}(g)) = \tau(g)$
« si $g \in$ Inertie ». La phrase « Mais la \ldots{} $\bar K$) le choix de »
est très cursive ; seuls quelques mots se lisent. Le vocabulaire des poids de
Hodge--Tate est postérieur.}

\pagerange{57}{59}
\section*{Trois feuilles de graphes (pages 57 à 59)}

Trois feuilles de papier quadrillé, au crayon, portent des lignes
polygonales issues de l'origine, des faisceaux de droites, et des
comparaisons de fractions. La première a pour en-tête « $n/17$ » et aligne
des petites lignes polygonales pour $n = 1, \ldots, 16$, en regard d'une
colonne marquée $\nu$ ; les inscriptions comparent des fractions de
dénominateurs 17, 23, 24 à des fractions de petits dénominateurs :
\[
\tfrac{5}{12} > \tfrac{7}{17}, \quad \tfrac{2}{7} < \tfrac{5}{17}, \quad
\tfrac{4}{17} < \tfrac{1}{4}, \quad \tfrac{4}{13} > \tfrac{7}{23}, \quad
\tfrac{2}{7} < \tfrac{7}{24}, \quad \tfrac{13}{17} < \tfrac{10}{13},
\quad \ldots
\]
Toutes celles qu'on lit en entier sont exactes. Les deux autres feuilles
reprennent la même chose avec des dénominateurs 15 et 16. Rien dans le
dossier ne dit à quoi servent ces graphes.\footnote{On peut penser à des
polygones de Newton, et à la question de savoir quels polygones à sommets
entiers approchent une droite de pente $n/17$ — question naturelle pour la
spécialisation des pentes d'un $F$-cristal, que Grothendieck étudie à cette
époque. C'est une conjecture de lecture, que les pages ne confirment pas. La
transcription donne toutes les inscriptions.}

\pagerange{60}{69}
\section*{Le groupe de Galois des $F$-cristaux filtrés (pages 60 à 68)}

\subsection*{Les données (pages 60 et 62)}

Soit $V$ un anneau de valuation discrète complet d'inégales
caractéristiques, de corps résiduel parfait $k$, de corps des fractions $L$ ;
$W = W(k) \subset V$, $K = W[1/p] \subset L$.\footnote{Dans cette section et
les suivantes, $K$ désigne, comme sur la page, le corps des fractions de
$W(k)$, et non plus le corps $p$-adique des pages 33 à 55.} Soit
$\mathscr{C}$ la catégorie des $F$-cristaux filtrés à isogénie près sur $V$ :
des couples $(E, F)$, $E$ un $K$-espace de dimension finie, $F$ un
endomorphisme $\sigma$-linéaire bijectif préservant un réseau, munis d'une
filtration de $E \otimes_{K} L$. C'est une catégorie tannakienne
$\mathbb{Q}_{p}$-linéaire, munie du foncteur fibre $\Delta$ « oubli » sur
$K$.\footnote{La page qualifie $\mathscr{C}$ de « ($\mathbb{Q}_{p}$-linéaire,
½ simple) », le « ½ simple » valant « semi-simple », et ajoute entre les
lignes « gr. déri. », de lecture incertaine. La catégorie des $F$-isocristaux
filtrés n'est pas semi-simple ; c'est une hypothèse de la page, ou une
restriction à une sous-catégorie, qu'on ne sait pas lire.}

La page dispose alors d'un « foncteur fondamental »
\[
T_{p}\colon \mathscr{C} \longrightarrow \operatorname{Mod}_{f}(\Gamma_{L},
\mathbb{Q}_{p}), \qquad \Gamma_{L} = \operatorname{Gal}(\bar L/L),
\]
vers les représentations $p$-adiques de dimension finie. C'est un foncteur
fibre sur $\mathbb{Q}_{p}$, et il identifie $\mathscr{C}$ à
$\operatorname{Rep}_{\mathbb{Q}_{p}}(G)$, $G$ un groupe proalgébrique sur
$\mathbb{Q}_{p}$.\footnote{L'existence d'un tel foncteur est précisément le
problème que Grothendieck pose à Nice en 1970 sous le nom de « foncteur
mystérieux ». La page l'utilise comme donné. Il a été construit par
Fontaine (anneau $B_{\mathrm{cris}}$, à partir de 1979), sur la
sous-catégorie des objets « faiblement admissibles », dont Colmez et Fontaine
ont montré (2000) qu'ils sont exactement ceux qui proviennent de
représentations. Le langage des groupes de Galois tannakiens est celui de la
thèse de Saavedra (1972). Le mot après « pro- » est biffé et récrit,
illisible les deux fois.}

Page 62 : on suppose $\mathscr{C}$ de type fini sur $\mathbb{Q}_{p}$. On a
un homomorphisme $\alpha\colon \Gamma_{L} \to G(\mathbb{Q}_{p})$, et $G$ est
l'enveloppe algébrique de $\alpha$ — l'adhérence de Zariski de son image —
si et seulement si $T_{p}$ est pleinement fidèle et son image essentielle
stable par sous-objets.\footnote{La page écrit $G(\mathbb{Q}_{\ell})$ ; il
s'agit de $G(\mathbb{Q}_{p})$. Elle énonce l'équivalence avec la seule pleine
fidélité (« dans ce cas seulement », lecture incertaine) ; il faut en plus la
stabilité par sous-objets. Elle ajoute « ou, ce qui revient au même, le
foncteur $\Delta$ » : on ne voit pas en quel sens, $\Delta$ étant un foncteur
fibre, et l'on ne reprend pas cette équivalence.}

\subsection*{Les cocaractères de Hodge--Tate et de Hodge (pages 62 et 64)}

Le composé
\[
\mathscr{C} \to \operatorname{Mod}_{f}(\Gamma_{L}, \mathbb{Q}_{p})
\xrightarrow{\ \otimes \mathbb{C}_{p}\ } \operatorname{Mod}_{f}(\Gamma_{L},
\mathbb{C}_{p}) \to \operatorname{Mod}_{f}(\mathbb{C}_{p})
\]
est bigradué — par le degré total et par la décomposition de Hodge(--Tate) —,
d'où deux cocaractères
\[
i_{1}, i_{2}\colon \mathbb{G}_{m,\mathbb{C}_{p}} \rightrightarrows
G_{\mathbb{C}_{p}}, \qquad i = i_{1}i_{2},
\]
$i$ étant le cocaractère central qui définit la graduation de
$\mathscr{C}$.\footnote{La graduation par le « degré total » suppose que
$\mathscr{C}$ est graduée (comme le sont les motifs, par le poids). Ce n'est
pas une conséquence de la définition de $\mathscr{C}$ donnée page 60 ; c'est
une hypothèse implicite de la page, qu'on rend explicite. La lecture de
l'égalité $i = i_{1}i_{2}$ est douteuse.} De même, le foncteur fibre
bigradué de Hodge $T_{\mathrm{Hdg}}$ (le gradué de la filtration) est défini
sur $L$ et correspond à un torseur $\mathcal{P}$ sous $G_{L}$, avec
$T_{\mathrm{Hdg}} = \mathcal{P} \wedge^{G} (T_{p} \otimes L)$, et à deux
cocaractères $j_{1}, j_{2}\colon \mathbb{G}_{m,L} \rightrightarrows G'_{L} =
\operatorname{ad}(\mathcal{P})$, de produit $i$.\footnote{La page dit
« $T_{\mathrm{Hdg}}$ sur $K$ » puis écrit un torseur sous $G_{L}$ ; la
filtration ne vivant que sur $E \otimes_{K} L$, c'est $L$ qui convient.}

La décomposition de Hodge--Tate s'écrit
\[
T_{p}(M) \otimes_{\mathbb{Q}_{p}} \mathbb{C}_{p} \simeq \bigoplus_{i,j}
T_{\mathrm{Hdg}}(M)^{i,j} \otimes \mathbb{C}_{p}(-i) .
\]
Tout isomorphisme $\rho\colon \mathbb{C}_{p}(1) \simeq \mathbb{C}_{p}$ en
déduit un isomorphisme bigradué $\tilde\rho\colon T_{p} \otimes \mathbb{C}_{p}
\simeq T_{\mathrm{Hdg}} \otimes \mathbb{C}_{p}$, c'est-à-dire un point
$\tilde\rho \in \mathcal{P}(\mathbb{C}_{p})$ ; si $\rho' = u\rho$, $u \in
\mathbb{C}_{p}^{*}$, alors $\tilde\rho'$ se déduit de $\tilde\rho$ par
$i_{1}(u^{-1})$. On obtient ainsi un morphisme d'espaces à opérations de
$\mathbb{C}_{p}^{*}$ et de $\Gamma_{L}$,
\[
\mathbb{C}_{p}(1)^{*} \longrightarrow \mathcal{P}(\mathbb{C}_{p}),
\]
où $\mathbb{C}_{p}^{*}$ opère sur $\mathcal{P}(\mathbb{C}_{p})$ via
$i_{1}^{-1}$ ; autrement dit un morphisme de
$\mathbb{G}_{m,\mathbb{C}_{p}}$-schémas de la droite épointée
$\check{\mathbb{V}}(\mathbb{C}_{p}(-1))^{*}$ dans $\mathcal{P}_{\mathbb{C}_{p}}$.
Il est compatible à $(i_{1}, i_{2})$ et $(j_{1}, j_{2})$ : $(j_{1}, j_{2})$ se
déduit par descente de $(i_{1}, i_{2})$ par l'isomorphisme
$G_{L}^{\mathcal{P}} \otimes_{L} \mathbb{C}_{p} \simeq G_{\mathbb{C}_{p}}$
associé à $\tilde\rho$.\footnote{La page écrit $\widetilde{u\rho} =
i_{1}(u^{-1})$, abréviation de ce qu'on a dit. Ce qu'elle construit est, en
termes d'aujourd'hui, un torseur de périodes de Hodge--Tate ; on n'en tire
aucune filiation.}

\subsection*{Le foncteur cristallin et son Frobenius (pages 66 et 68)}

Il y a aussi un foncteur fibre $T_{\mathrm{cr}}$ sur $K$ — l'isocristal
sous-jacent —, d'où un torseur $Q$ sous $G_{K}$ avec $T_{\mathrm{cr}} \simeq
Q \wedge^{G_{K}} (T_{p})_{K}$. Le foncteur $T_{\mathrm{cr}} \otimes_{K} L$ est
filtré, de gradué $T_{\mathrm{Hdg}}$ ; d'où un torseur $R$ sous un sous-groupe
unipotent $U' \subset G' = G_{L}^{\mathcal{P}}$,
\[
R \hookrightarrow \operatorname{Isom}(Q_{L}, \mathcal{P}) .
\]
Enfin le Frobenius $F\colon T_{\mathrm{cr}}^{\sigma} \xrightarrow{\sim}
T_{\mathrm{cr}}$ correspond à un isomorphisme
\[
Q^{\sigma} \simeq Q ,
\]
ce que la page 68 encadre, et ce que le griffonnage de la page 72 reprend.\footnote{La page appelle ici $T_{\mathrm{DR}}$ le foncteur qu'elle appelait
$T_{\mathrm{cr}}$ quelques lignes plus haut ; l'indice est surchargé. Un bloc
encadré et hachuré introduisait un cocaractère $k_{2}\colon
\mathbb{G}_{m,K} \to G'' = \operatorname{ad}(Q)$ et un unipotent $U''$ ; il
est barré, et seul $G'' \supset U''$ en survit dans l'encadré suivant. Un
passage barré de croisillons se lit « il faudra expliquer \ldots{} cela
correspond \ldots{} $i_{2}$ \ldots{} structures bigraduées ».}

\subsection*{Un second tapuscrit (pages 61 à 73)}

Six feuillets d'un autre tapuscrit (numérotés 28, 29, 31, 33, 34 et 35)
s'intercalent dans ces pages : $S$-couples lisses $(Y, X)$ de codimension
$c$, suite de Kummer et morphisme $\varepsilon_{t}$ de (3.9), Théorème
(3.10) qui donne, pour chaque couple lisse, un isomorphisme canonique
$\varphi(Z, Y)$, son diagramme de transitivité (3.10.1) et sa démonstration
par récurrence sur la codimension, puis un énoncé de changement de base pour
$(R^{q}j_{*})g^{*}F$ démontré par la suite spectrale de Hochschild--Serre.
Les interventions à la main sont des lettres grecques et des renvois
— « (IX 3.2) », « (VIII 5.5) », « (cf. XII 6.5) » — et, en marge d'une
démonstration, « démonstration peu naturelle ».\footnote{La lettre $\varphi$
de ces feuillets est celle du tapuscrit, sans rapport avec le $\varphi$ du
dossier. Leur numérotation entre en collision avec celle du tapuscrit des
pages 6 à 14 : c'est au moins une autre partie, sans doute un autre texte.
On ne l'identifie pas.}

\pagerange{70}{74}
\section*{Les structures d'un motif (pages 70 à 74)}

\subsection*{Sur un ouvert de $\operatorname{Spec}\mathbb{Z}$ (page 70)}

La page 70 dresse la même liste pour la catégorie $\mathscr{C}$ des
(iso)motifs sur un ouvert $S \subset \operatorname{Spec}\mathbb{Z}$ :

\begin{itemize}
\item le foncteur de Betti $T_{B}$, sur $\mathbb{Q}$, relatif à $\mathbb{C}$,
et son groupe $G$ sur $\mathbb{Q}$ ;
\item les réalisations $\ell$-adiques $T_{\ell} \simeq T_{B} \otimes
\mathbb{Q}_{\ell}$, d'où $\Gamma_{\mathbb{Q}} \to G(\mathbb{Q}_{p})$, non
ramifié aux nombres premiers de $S$ autres que $p$ ;
\item pour chaque place $p \in S$, finie ou infinie, le groupe de
décomposition $\Gamma_{\mathbb{Q}_{p}}$ (et $\mathbb{Z}/2\mathbb{Z}$ pour
$p = \infty$), et des paires de cocaractères $i_{1}(\infty), i_{2}(\infty)$
de $G_{\mathbb{C}}$, $i_{1}(p), i_{2}(p)$ de $G_{\mathbb{C}_{p}}$ ;
\item le foncteur de De Rham $T_{\mathrm{DR}}$ sur $\mathbb{Q}$, de groupe
$G' = \operatorname{ad}(\mathcal{P})$ pour un torseur $\mathcal{P}$ sous
$G$, filtré, c'est-à-dire muni d'une classe de cocaractères
$j_{2}\colon \mathbb{G}_{m,\mathbb{Q}} \to G'$ définie à conjugaison près par
le radical unipotent $U' \subset G'$ (soit un torseur $R$ sous $U'$) ;
\item des Frobenius $F_{p} \in G'(\mathbb{Q}_{p})$ pour $p \in S$ fini, et
$F_{\infty} \in G'(\mathbb{R})$ d'ordre 2 ;
\item un élément $\rho^{(\infty)} \in \mathcal{P}(\mathbb{C})$ qui transforme
la classe des $j_{2}$ en celle des $i_{2}^{(\infty)}$ — l'isomorphisme de
comparaison de Betti à De Rham ;
\item pour chaque $p$, une application
$\mathbb{C}_{p}(1)^{*} \times U'(\mathbb{C}_{p}) \to
\mathcal{P}(\mathbb{C}_{p})$, $(\lambda, j_{2}) \mapsto \rho^{(p)}(j_{2},
\lambda)$, qui transforme $j_{2}$ en $i_{2}^{(p)}$, compatible aux
opérations de $\mathbb{C}_{p}^{*} \times U'(\mathbb{C}_{p})$ et de
$\Gamma_{\mathbb{Q}_{p}}$ — la comparaison $p$-adique.
\end{itemize}

Aucune de ces structures n'est construite sur la page : c'est une liste de
ce qu'un motif doit porter.\footnote{La page note « (iso) » au-dessus de « motifs », « ½ simples », et
dans la marge « relations de $T_{B}$ et $T_{\mathrm{DR}}$ ». Elle écrit
$C^{p}$ ou $C^{(p)}$ pour $\mathbb{C}_{p}$. La comparaison $p$-adique pour les
variétés à bonne réduction, qui donne ici $\rho^{(p)}$, a été démontrée par
Fontaine--Messing et Faltings dans les années 1980 ; sur la page elle est une
structure postulée.}

\subsection*{Au-dessus d'une base (page 72)}

Le haut de la page 72, écrit en travers, est un griffonnage où se lisent
$H^{*}_{B}$, $H_{\mathrm{Hdg}}$, $H_{\mathrm{DR}}$, $H_{\mathrm{cris}}$,
« $Q^{\sigma} \simeq Q$ » et un petit tableau « $T_{B}$ : $G_{B}$ ;
$T_{\mathrm{DR}}$ : $G_{\mathrm{DR}}$ ; $T_{\mathrm{Hdg}}$ ». Le bas énumère
les structures d'un motif au-dessus d'une base $S$ :

\begin{enumerate}
\item[a)] un module $M$ à connexion intégrable sur $S$, avec filtration ;
\item[b)] un Frobenius $F_{p}\colon M(p)^{\sigma} \to M(p)$ (théorie des
cristaux) ;
\item[c)] $M(\mathbb{R})$ analytique filtré sur $S(\mathbb{R})$, avec un
automorphisme $F_{\infty}$ d'ordre 2 ;
\item[d)] un système local de Betti $T$ sur $S(\mathbb{C})^{\mathrm{an}}$,
soit une représentation de $\pi_{1}(S(\mathbb{C})^{\mathrm{an}}, s)$ ;
\item[e)] un système local $p$-adique $T_{p}$ sur $S[1/p]$, avec
$T_{p} \simeq T \otimes \mathbb{Z}_{p}$ sur $S_{\mathbb{C}}$ ;
\item[f$_{\infty}$)] la comparaison $M^{\mathrm{an}}_{\mathbb{C}} \simeq T
\otimes_{\mathbb{Z}} \mathcal{O}^{\mathrm{an}}$, compatible aux connexions ;
\item[f$_{p}$)] laissée en blanc.
\end{enumerate}

La comparaison $p$-adique est la seule rubrique sans contenu.\footnote{La page écrit f$_{\infty}$) « $M^{\mathrm{an}}_{\mathbb{C}} \simeq T
\otimes_{\mathbb{Z}} \mathbb{C}$ » : au sens propre, c'est l'isomorphisme
entre les sections horizontales de $M^{\mathrm{an}}$ et $T \otimes
\mathbb{C}$, ou, de façon équivalente, $M^{\mathrm{an}} \simeq T \otimes
\mathcal{O}^{\mathrm{an}}$ avec la connexion triviale sur le second membre.
La place de la comparaison $p$-adique, f$_{p}$), est marquée par quatre
tirets.}

\subsection*{Les conjectures (page 74)}

Une liste, avec en regard la structure que chacune concerne :

\begin{itemize}
\item conjecture de Hodge (généralisée) — la décomposition de Hodge sur
$\mathbb{C}$ ;
\item conjecture de Tate (généralisée) — l'action de Galois, les Frobenius ;
\item conjecture des Frobenius en cohomologie de De Rham (à connexion) — le
Frobenius cristallin ;
\item conjecture du double réseau, précédée d'un « ? » et suivie de
« (revoir) » ;
\item « Tate groups » (généralisés), sur les schémas de type fini sur
$\mathbb{Z}$ ou sur un corps, et sur les anneaux locaux d'inégales
caractéristiques.
\end{itemize}

Aucune n'est discutée plus avant dans le dossier.\footnote{Un mot illisible précède chaque « généralisée ». Ce que la page
entend par « double réseau » n'est pas dit ; la question des deux réseaux de
la page 49 en est un candidat, que rien ne confirme. Les « Tate groups » sont
peut-être l'analogue, pour l'image de Galois, du groupe de Mumford--Tate ;
c'est aussi une conjecture de lecture.}

\pagerange{77}{80}
\section*{Le complexe co-Lie des groupes de type multiplicatif (pages 77 à 80)}

\subsection*{Le calcul (pages 77 et 78)}

Soit $G = D(M)$ le groupe de type multiplicatif dual d'un groupe constant
tordu $M$ sur $S$, de type fini. On résout $M$ par la résolution plate
canonique à deux crans
\[
(1.1) \qquad 0 \to R \xrightarrow{\ v\ } \mathbb{Z}[M] \xrightarrow{\ u\ } M
\to 0, \qquad u([m]) = m,
\]
d'où, par dualité, une résolution de $G$ par des tores
\[
(1.2) \qquad 0 \to G \to T \to T' \to 0, \qquad T = D(\mathbb{Z}[M]) \simeq
\mathbb{G}_{m}^{M}, \quad T' = D(R),
\]
de dimension relative infinie si $M$ est infini. Le triangle de
transitivité appliqué à $G = \operatorname{Ker}(T \to T')$, $T$ et $T'$ lisses,
donne
\[
\ell^{G} \simeq [\omega_{T'} \to \omega_{T}] \simeq [R \otimes_{\mathbb{Z}}
\mathcal{O}_{S} \to \mathbb{Z}[M] \otimes_{\mathbb{Z}} \mathcal{O}_{S}],
\qquad\text{soit}\qquad
(1.3) \quad \ell^{G} \simeq M \otimes^{L}_{\mathbb{Z}} \mathcal{O}_{S} .
\]
Si $\mathcal{O}_{S}$ est une $\Lambda$-algèbre (par exemple $\Lambda =
\mathbb{Q}$ ou $\mathbb{Z}/n\mathbb{Z}$),
\[
(1.4) \qquad \ell^{G} \simeq (M \otimes^{L}_{\mathbb{Z}} \Lambda)
\otimes^{L}_{\Lambda} \mathcal{O}_{S},
\]
et si $\Lambda$ est un corps, par exemple $\mathbb{F}_{p}$, le second produit
tensoriel n'a pas besoin d'être dérivé :
\[
(1.4') \qquad \ell^{G} \simeq (M \otimes^{L}_{\mathbb{Z}} \mathbb{F}_{p})
\otimes_{\mathbb{F}_{p}} \mathcal{O}_{S} .
\]

\emph{1.5. Le cas d'une descente.} Supposons que $M$ provienne, par un
revêtement principal $S'$ de $S$ de groupe $\Pi$ discret ou profini, d'une
représentation (continue) de $\Pi$ dans un groupe abélien $M_{0}$ — par
exemple $M$ isotrivial, à fibres de type fini, $S$ normal et connexe, $S'$
le revêtement universel. On a un morphisme de topos
\[
(1.6) \quad f\colon S \longrightarrow B_{\Pi}, \qquad (1.7) \quad M =
f^{-1}(M_{0}),
\]
pour la topologie étale, fppf ou fpqc, et, en annelant $S$ et $B_{\Pi}$ par
$\Lambda$, $M \otimes^{L}_{\mathbb{Z}} \Lambda = f^{*}(M_{0}
\otimes^{L}_{\mathbb{Z}} \Lambda)$, d'où
\[
(1.8) \qquad \ell^{G} \simeq f^{*}(M_{0} \otimes^{L}_{\mathbb{Z}} \Lambda)
\otimes^{L}_{\Lambda} \mathcal{O}_{S} ,
\]
le $L$ du second produit tensoriel étant inutile si $\Lambda$ est un
corps.\footnote{On note $\Pi$ le groupe du revêtement, que la page appelle $\pi$.
Le « 1.5 » est écrit sur un « 2. ». Un numéro de formule entre (1.7) et
(1.8) est surchargé et illisible. La page attribue le triangle de
transitivité à Mazur et Roberts ; voir la note sur le complexe co-Lie dans
les conventions.}

\subsection*{La classe $\xi$ (pages 79 et 80)}

De (1.3) on tire
\[
(2.1) \qquad \omega_{G} = M \otimes_{\mathbb{Z}} \mathcal{O}_{S}, \qquad
n_{G} = \operatorname{Tor}_{1}^{\mathbb{Z}}(M, \mathcal{O}_{S}),
\]
et si $S$ est de caractéristique $p$,
\[
(2.2) \qquad \omega_{G} = (M/pM) \otimes_{\mathbb{F}_{p}} \mathcal{O}_{S},
\qquad n_{G} \simeq {}_{p}M \otimes_{\mathbb{F}_{p}} \mathcal{O}_{S} ;
\]
si $M$ est tué par $p$, c'est-à-dire $G$ tué par $p$,
\[
(2.2') \qquad \omega_{G} \simeq n_{G} \simeq M \otimes_{\mathbb{F}_{p}}
\mathcal{O}_{S} .
\]
C'est le phénomène $n = \omega$ de la page 12, rencontré cette fois pour un
groupe de type multiplicatif, et calculé directement.

Le complexe $\ell^{G}$, à deux faisceaux de cohomologie, porte une classe
canonique — son invariant de Postnikov —
\[
\xi \in \operatorname{Ext}^{2}_{\mathcal{O}_{S}}\bigl(M \otimes_{\mathbb{Z}}
\mathcal{O}_{S},\ \operatorname{Tor}_{1}^{\mathbb{Z}}(M, \mathcal{O}_{S})\bigr),
\]
qui provient, en présence de $\Lambda$, d'une classe
$\xi_{\Lambda} \in \operatorname{Ext}^{2}_{\Lambda}(M \otimes \Lambda,
\operatorname{Tor}_{1}^{\mathbb{Z}}(M, \Lambda))$. Si $\Lambda =
\mathbb{Z}/n\mathbb{Z}$ et $M$ est un $\Lambda$-module localement libre de
type fini, $M \otimes \Lambda \simeq M$ et
$\operatorname{Tor}_{1}^{\mathbb{Z}}(M, \Lambda) = {}_{n}M \simeq M$, donc
\[
\xi_{\Lambda} \in H^{2}(S, \mathcal{E}nd(M)), \qquad
\xi \in H^{2}(S, \mathcal{E}nd(E)) \quad (E = M \otimes_{\Lambda}
\mathcal{O}_{S}),
\]
$\xi$ étant l'image de $\xi_{\Lambda}$ par la flèche canonique
$H^{2}(S, \mathcal{E}nd(M)) \to H^{2}(S, \mathcal{E}nd(M) \otimes_{\Lambda}
\mathcal{O}_{S})$.\footnote{Trois mots illisibles suivent « Alors » dans la
phrase qui donne ces identifications ; on les a rétablies par le calcul.
Une hypothèse de platitude de $\mathcal{O}_{S}$ sur $\Lambda$ est écrite
puis biffée.} Sous les conditions de 1.5, $M_{0}$ étant plat sur $\Lambda$,
$\xi_{\Lambda}$ est l'image d'une classe
\[
\xi_{0} \in H^{2}(B_{\Pi}, \mathcal{E}nd(M_{0})) = H^{2}(\Pi,
\operatorname{End}(M_{0})) .
\]
Pour la calculer, on se ramène par fonctorialité au cas universel
$M_{0} = (\mathbb{Z}/n\mathbb{Z})^{d}$, $\Pi = \mathrm{GL}(d,
\mathbb{Z}/n\mathbb{Z})$. D'après Illusie, c'est l'obstruction à relever la
représentation de $\Pi$ dans $M_{0}$ en une représentation dans un module
libre sur $\mathbb{Z}/p^{2}\mathbb{Z}$ — pour $n = p$, la classe de
l'extension $1 \to M_{d}(\mathbb{F}_{p}) \to \mathrm{GL}(d,
\mathbb{Z}/p^{2}) \to \mathrm{GL}(d, \mathbb{F}_{p}) \to 1$. Elle est nulle
pour $d \le 1$, puisque $\mathbb{F}_{p}^{*}$ se relève par les
représentants de Teichmüller. La page ajoute qu'elle n'est « sans doute pas
nulle » pour $d \ge 2$ ; c'est vrai en général, mais pas sans
exception.\footnote{Pour $p = 2$ et $d = 2$, $\mathrm{GL}(2, \mathbb{F}_{2})
\simeq \mathfrak{S}_{3}$ se relève à $\mathrm{GL}(2, \mathbb{Z})$ par sa
représentation de permutation réduite, donc à $\mathrm{GL}(2,
\mathbb{Z}/4)$ : l'extension est scindée et $\xi_{0} = 0$. Il y a d'autres
petits cas scindés ; on ne les énumère pas.}

Enfin, si $S'$ est le revêtement universel de $S$, l'application
$H^{2}(B_{\Pi}, -) \to H^{2}(S, -)$ est injective — dans la suite spectrale de
Hochschild--Serre du revêtement, le terme $H^{0}(\Pi, H^{1}(S', -))$ qui
pourrait tuer $H^{2}(\Pi, -)$ est nul, $S'$ étant simplement connexe —, de
sorte que la classe $\xi_{\Lambda} \in H^{2}(S, \mathcal{E}nd(M))$ est non
nulle dès que $\xi_{0}$ l'est. Le dossier s'arrête là, sur un énoncé qui
rejoint sa première page : le complexe co-Lie d'un groupe tué par $p$ ne se
réduit pas à ses deux faisceaux de cohomologie, et ce qui l'en empêche est
une obstruction de relèvement modulo $p^{2}$.\footnote{L'injectivité est
affirmée par la page ; la justification par Hochschild--Serre est la
nôtre, et vaut pour les coefficients localement constants provenant de
$B_{\Pi}$.}

\end{document}
